\documentclass[11pt]{article}

\usepackage[T1]{fontenc}
\usepackage[utf8]{inputenc}
\usepackage{lmodern}
\usepackage[margin=1in]{geometry}
\usepackage{amsmath,amssymb,amsthm,mathtools}
\usepackage{enumitem}
\usepackage{booktabs}
\usepackage{array}
\usepackage{longtable}
\usepackage{xcolor}
\usepackage{microtype}
\usepackage{etoolbox}
\usepackage{url}
\usepackage[colorlinks=true,linkcolor=blue!50!black,citecolor=green!35!black,urlcolor=blue!60!black]{hyperref}
\usepackage[capitalize,nameinlink]{cleveref}

% ---------------------------------------------------------------------------
% Theorem environments
% ---------------------------------------------------------------------------
\theoremstyle{plain}
\newtheorem{theorem}{Theorem}[section]
\newtheorem{lemma}[theorem]{Lemma}
\newtheorem{corollary}[theorem]{Corollary}
\newtheorem{proposition}[theorem]{Proposition}
\newtheorem{conjecture}[theorem]{Conjecture}
\newtheorem{question}[theorem]{Question}
\newtheorem{problem}[theorem]{Problem}
\theoremstyle{definition}
\newtheorem{definition}[theorem]{Definition}
\newtheorem{remark}[theorem]{Remark}
\newtheorem{example}[theorem]{Example}

% The Lean proof is the proof; the text gives a faithful outline of it.
\newenvironment{outline}[1][Outline of the formal proof]{\begin{proof}[#1]}{\end{proof}}
% Results that are not formalized get complete proofs, labelled as such.
\newenvironment{informalproof}{\begin{proof}[Proof (not formalized)]}{\end{proof}}

% ---------------------------------------------------------------------------
% Lean references
%   \lean{name}            a declaration, relative to the namespace SwapRigidity
%   \formal{file}{names}   the formal counterpart of a numbered result: a file relative to
%                          lean/SwapRigidity/ (or, when every declaration belongs to an earlier
%                          library, a file of that library) and a comma-separated list of declarations;
%                          names from earlier libraries carry their namespace prefix
%   \partformal{file}{names}{gap}
%                          a numbered result whose formal counterpart covers all but the
%                          part described in {gap}; the gap is proved or cited in the text
%   \notformalized         a numbered result that has no formal counterpart
%   \cited{where}          a result of an earlier paper used without proof: where it is proved,
%                          and whether it is Lean-verified there
% ---------------------------------------------------------------------------
% Lean names are typeset like paths, with line breaks allowed after "_", "." and "/".
\DeclareUrlCommand\leanname{\def\UrlFont{\ttfamily\upshape}}
\renewcommand{\UrlBreaks}{\do\_\do\.\do\/\do\-}
\newcommand{\lean}[1]{\leanname{#1}}
\newcommand{\leanfile}[1]{\leanname{#1}}
\newcommand{\leanlist}[1]{\begingroup\def\leansep{}%
  \renewcommand*{\do}[1]{\leansep\leanname{##1}\def\leansep{, }}\docsvlist{#1}\endgroup}
\newcommand{\formal}[2]{\par\noindent{\small\upshape\raggedright\textsf{Lean:} \leanlist{#2}\unskip
  \nobreak\hfil\penalty50\hskip1em\hbox{}\nobreak\hfill[\leanfile{#1}]\parfillskip=0pt\par}%
  \smallskip\noindent\ignorespaces}
\newcommand{\partformal}[3]{\par\noindent{\small\upshape\raggedright\textsf{Lean (partial):} \leanlist{#2}\unskip
  \nobreak\hfil\penalty50\hskip1em\hbox{}\nobreak\hfill[\leanfile{#1}]\parfillskip=0pt\par}%
  \par\noindent{\small\upshape\textsf{Not formalized:} #3}\par\smallskip\noindent\ignorespaces}
\newcommand{\notformalized}{\par\noindent{\small\upshape\textsf{Not formalized;} a complete proof is given
  below.}\par\smallskip\noindent\ignorespaces}
\newcommand{\cited}[1]{\par\noindent{\small\upshape\textsf{Cited:} #1}\par\smallskip\noindent\ignorespaces}

% ---------------------------------------------------------------------------
% Notation
% ---------------------------------------------------------------------------
\newcommand{\N}{\mathbb{N}}
\newcommand{\Z}{\mathbb{Z}}
\newcommand{\Q}{\mathbb{Q}}
\newcommand{\R}{\mathbb{R}}
\newcommand{\C}{\mathbb{C}}
\newcommand{\F}{\mathbb{F}}
\newcommand{\T}{\mathsf{T}}
\newcommand{\clog}[1]{\lceil \log_2 #1 \rceil}
\newcommand{\abs}[1]{\lvert #1 \rvert}
\newcommand{\card}[1]{\lvert #1 \rvert}
\newcommand{\eps}{\varepsilon}
\newcommand{\prev}{\mathrm{prev}}
\newcommand{\cost}{\mathrm{cost}}
\newcommand{\gn}{\mathrm{gn}}
\newcommand{\Sh}{\mathrm{Sh}}
\newcommand{\cl}{\mathrm{cl}}
\newcommand{\Aut}{\mathrm{Aut}}
\newcommand{\minrk}{\mathrm{minrk}}
\newcommand{\MAIS}{\mathrm{MAIS}}
\newcommand{\Crux}{\mathrm{Crux}}
\newcommand{\TCP}{\mathrm{TCP}}
\DeclareMathOperator{\rank}{rank}
\DeclareMathOperator{\dist}{dist}
\DeclareMathOperator{\poly}{poly}

% ---------------------------------------------------------------------------
% Notation of this paper
% ---------------------------------------------------------------------------
\newcommand{\sw}{\tau}
\newcommand{\X}{\mathcal{X}}
\newcommand{\Hk}{\mathcal{H}}
\newcommand{\cA}{\mathcal{A}}
\newcommand{\tr}{\mathsf{T}}
\DeclareMathOperator{\im}{im}
\DeclareMathOperator{\spn}{span}
\DeclareMathOperator{\chr}{char}
\newcommand{\Mat}{\mathrm{M}}


\title{Order, Symmetry and the Swap:\\ A Synthesis of Matrix Transposition on Tapes}
\author{Hans Brende \and Claude Fable 5.2 (Anthropic)}
\date{November 2026}

\begin{document}

\maketitle

\begin{abstract}
Can a deterministic multitape Turing machine transpose a $d\times d$ matrix with $\lceil\log_2d\rceil$-bit entries in
time linear in the input length~$L$? Sixteen earlier papers of this series study the question without deciding it. This
paper is a synthesis of them, a diagnosis of where they stop, and two new theorems.

The synthesis states the main results of the series in one place, with their main hypotheses. The diagnosis is
this. For oblivious machines the question is now one about a single function, the best rate of a storage code on the
bounded-depth cone relation that is left of a linear-time schedule after a genie has revealed $\eps L$ bits: a
transposer needs rate $1-\eps$ for every~$\eps$, pairing gives rate $\tfrac12$, and for the relations of real machines
in general no rate above $\tfrac12$ is achieved and no useful bound below one is proved. Conditions on head motion that
ignore the data are met by a machine that transposes nothing. Bounds for codes that ignore the machine fail, because
abstract relations with the two properties used so far, bounded degree and paired-distance separation, carry codes of
rate tending to one. Counting passes or cuts saturates one scale at a time. The common reason is that the $\log d$ of
the folklore algorithm is a sum over address scales, and every tool that applies to all machines looks at one scale or
at none; the logarithm is proved only where either the coding or the geometry is tied down.

The theorems concern the swap $\sw$ on $V\otimes V$, $V=F^d$, and the code $\ker(\sw-K)$ of matrices that an operator
$K=\sum_aA_a\otimes B_a$ transposes correctly. \emph{Order}: for every complete flag of $V$, the code has dimension at most
the sum of the dimensions of the images of the flag under the $B_a$, whatever the $A_a$ are; so two terms, commuting
families, and, in characteristic $p$ with $d=p^m$, all affine index maps stay at the dimension $d(d+1)/2$ of the
symmetric matrices. \emph{Symmetry}: if each row of a matrix is a fixed linear function of the matrix seen from that
row's position, the same function on each of $M$ intervals of rows, then the matrix ranges over a space of dimension at
most $(d^2+Md)/2$, on a cycle and on a line, with equality for equal intervals, where the bound is the dimension of the
coarse $M\times M$ block transposition. The proof intersects the code with its own translate. It works in every
characteristic and uses a property of machines that cone relations do not record: translated outputs share their
decoder. Variants price each change of the decoder of a single output entry, from one row to the next, at half a
symbol, where a change of class costs half a row, also in the presence of advice over a finite field; bound decoders
with few modulation frequencies; and show that every address digit must be treated asymmetrically, which one unshuffle
does for all digits at once. For classes of rows that no single shift respects the bound
can fail: at $d=6$ two such classes carry a code of dimension $25>24$, and, by a rule taken from a memo to this paper,
two classes carry codes that exceed the bound by $d/4-1$ at every $d=4k\ge8$, over every field; for every fixed number
$M\ge2$ of classes the gain over the bound is unbounded. A later paper of the series, cited in this revision, bounds
two classes of any shape by $\tfrac{41}{64}d^2+O(d^{3/2})$ and shows that the second theorem needs the code, but not
the decoders, to be linear; for sets of matrices and arbitrary decoders two bounds of this paper remain, the
elimination step of the first theorem and the bound for rules that copy one entry, as instances of the
acyclic-subgraph bound. The same bound shows that a set of matrices over $q$ symbols in which every entry is a
function of at most two entries, at positions that move with the row, has at most $q^{d(d+1)/2+d-1}$ elements, and
that reading any number of entries at boundedly many offsets gives no more than $q^{(1/2+o(1))d^2}$; whether three
entries at arbitrary positions of this kind can give more is open.

Both theorems are statements about a linear one-hop model, not about machines, and the question remains open. The
coordinate form of the first theorem, its attainment, the pigeonhole lemma, the second theorem on a cycle for arbitrary
shift sets and on a line, its attainment for equal intervals, its form for changes at single entries with and
without advice, and the twist-class variant are formalized in Lean~4 on the standard axioms. The corollaries that need an algebraically closed
field or the structure of $p$-groups, the copy-rule bound, the modulation and digit corollaries, the passage from
flags to bases, one count in a sharper form of the bound for single entries, the gains for a fixed number of
classes and the bounds for sets are proved in the text and not formalized. Exact computations, at sizes up to six and for the two-class rule and its blow-ups at larger sizes, are
reported as numerical remarks.

\end{abstract}

\begin{quote}\small
\textbf{Contributions.} Hans Brende chose the problem, fixed the verification standard, directed the research, and asked for
this synthesis. Claude Fable~5.2, a large language model developed by Anthropic, read the sixteen earlier papers and their
memos, wrote the synthesis and the diagnosis, developed the mathematics of \cref{sec:order,sec:symmetry} and the exact
computations of \cref{sec:small}, directed the Lean formalization, and wrote this paper. Independent instances of the same
model refereed the two theorems; one of them contributed \cref{cor:affine}, the line form of \cref{thm:TP}
(\cref{thm:line}) and the counterexample of \cref{rem:line-fails}. The revision that follows the
memo~\cite{ClassUniformMemo2026} was written by Claude Fable~5.2 at Hans Brende's request, with the Lean proofs of
\cref{prop:entry,cor:entry-advice}; a further instance, which refereed it, pointed out that the charge for revealed
bits confused two units, and contributed the bound of \cref{prop:entry} with its proof and the combination with advice
that became \cref{cor:entry-advice}. The rule of \cref{prop:two-class} is the memo's rule with a general
coefficient, and the proof of its dimension follows the memo's, where Gemini~4 Argon found the rule and part of the
rank count and Claude Opus~5.5 completed the proof; Claude Fable~5.2 observed that the rule works with a general
coefficient and over every field, and adapted the proof. A further instance, which refereed the
revision that follows the memo's update, observed that a blow-up with a code for each position settles every fixed
number of classes, and gave the argument of \cref{lem:mixed-blowup,cor:every-M}. The revision that follows the two
later papers~\cite{AdaptiveTransfer2026,SharedDecoders2026}, by Hans Brende and GPT-6 Astra, was written by Claude
Fable~5.2, which read both papers, rebuilt their Lean libraries and wrote \cref{sec:sets} and the new form of
\cref{prob:machines}. A further instance refereed it against the two papers. It pointed out that
\cref{prop:sets} is an instance of the acyclic-subgraph bound, gave its form with rows and columns exchanged, and
observed that the examples of the second paper lack only a bounded window, so that the question with one function
for all rows had to be asked as well. In a second pass it contributed the marker rows of \cref{sec:sets} and their
consequence for \cref{prob:machines}(b). The revision after that, on \cref{prob:machines}(b), is by Claude
Fable~5.2: it found \cref{prop:order,prop:mutual}, with the ranking lemma and the two counting arguments, and ran
the searches reported at the end of \cref{sec:sets}. A further instance refereed the proofs and the scripts and found no error in the proofs. It contributed the
examples after \cref{prop:order} (that the ranking lemma is sharp on cycles, the relation at $d=3$ that attains the
bound of window two, and the lower bound~$1$ for the constant at even~$d$), observed that parts (c) and~(d) with a
multiplier prime to~$d$ follow from \cref{prop:sets} after a renumbering, and that one entry of the table at the end
of \cref{sec:sets} can be turned into a proof. In a second pass it read the text of this revision against the
seventh paper and corrected two sentences of the commentary on \cref{prob:machines}.
\end{quote}

\tableofcontents
\clearpage

\section{Introduction}\label{sec:intro}

\subsection{The problem}

Can a deterministic multitape Turing machine transpose a $d\times d$ matrix with $\clog d$-bit entries in time linear
in the input length $L=d^2\clog d$? We call the affirmative statement~(T); \cref{sec:question} states it precisely. The
folklore algorithm takes $O(L\log d)$ steps. Transposition is the data-movement step of fast integer multiplication on
tapes, and by a reduction of Harvey and van der Hoeven~\cite{HvdH2025} a linear-time multiplier would give~(T). So a
proof that (T) fails would be the first superlinear lower bound for off-line multiplication on multitape machines, and a
proof that it holds would be a new way of moving data.

Sixteen papers and eight memos of this series have studied the question since October~2026. They contain reductions,
restricted impossibility theorems, constructions of codes and many formal proofs. They do not contain the answer, and a
reader who starts today has to read all of them to learn where the difficulty is. This paper is written for that reader.

\subsection{What this paper does}

It does three things.

\paragraph{A synthesis (\cref{sec:series}).} We state the main results of the series in one place, by theme: the
reductions, what head motion must do, depth reduction, what is proved about storage codes, what is false about them, and
adaptive machines. Nothing is reproved. Every statement is cited, with its main hypotheses.

\paragraph{A diagnosis (\cref{sec:wall}).} For oblivious machines the series has reduced~(T) to one function: the best
rate $r(\eps)$ of a storage code on the bounded-depth cone relations that remain of linear-time schedules when a genie
reveals $\eps L$ bits. A transposer needs $r(\eps)\ge1-\eps$ for every~$\eps$. The matching of each stored symbol with
one partner gives $r\ge\tfrac12$, and whether $r$ exceeds $\tfrac12$ at all on the relations of real machines is open in
both directions. Three kinds of argument have reached known limits:
\begin{itemize}[leftmargin=1.5em]
\item conditions on motion that ignore the data, because a machine that transposes nothing satisfies the cut
inequalities behind them; conditions that do see the outputs are single-scale, and one cut serves every pair of row
sets;
\item bounds for codes that ignore the machine, because abstract relations with the two properties used so far,
bounded degree and paired-distance separation, carry codes of rate tending to one;
\item counting passes, because it saturates one scale at a time.
\end{itemize}
They stop for one reason. The logarithm of the folklore bound is a sum over $\log_2d$ address scales, each costing $L$,
and every tool of the series that applies to all machines looks at one scale or at none; the logarithm is proved only
where the coding or the geometry is tied down. The missing statement has to tie what is known to where it can be, for
all schedules, and the coding properties proved so far are properties of the cone \emph{relation}, which forgets that
neighbouring outputs of a machine are produced by the same rule.

\paragraph{Two theorems (\cref{sec:order,sec:symmetry}).} We prove two rigidity theorems for the swap in a model where
the code and its decoder are visible together. Let $V=F^d$ over any field, let $\sw$ be the swap on $V\otimes V$, that
is, $X\mapsto X^{\tr}$ on matrices, and let $K=\sum_aA_a\otimes B_a$ act by $X\mapsto\sum_aA_aXB_a^{\tr}$. The code
$\ker(\sw-K)$ is the set of matrices that $K$ transposes correctly. The symmetric matrices, for $K=1$, have dimension
$d(d+1)/2$.

\begin{itemize}[leftmargin=1.5em]
\item \emph{Order (Theorem~FL, \cref{thm:FL}).} For every complete flag $W_1\subset\dots\subset W_d$ of $V$, with
$g(i)=\dim\sum_aB_a(W_i)$,
\[
  \dim\ker(\sw-K)\le\sum_{i=1}^dg(i),
\]
whatever the $A_a$ are. So two Kronecker terms never beat the symmetric matrices, over any field; neither do commuting
$B_a$, for instance any convolutions of any number of rows; and neither does any combination of rotations, reflections,
dilations and decimations when $d=p^m$ in characteristic~$p$. The bound is attained for every profile~$g$.
\item \emph{Symmetry (Theorem~TP, \cref{thm:TP}).} For any three operators on an $N$-dimensional space,
$\dim\ker(\sw-K)\le\frac12(N+\dim\ker[\sw,T]+\rank[K,T])$, an elementary lemma. For the swap, $\dim\ker[\sw,T]\ge d$
always. If row $c$ of $X$ is a fixed linear function of $X$ with its columns rotated by $c$, the same function for all
$c$ in each of $M$ intervals, then
\[
  \dim\{X\}\le\frac{d^2+Md}2 ,
\]
on a cycle and on a line, with equality for equal intervals: there the largest code has the dimension of the coarse
$M\times M$ block transposition. For $M$ equal intervals flags give only rate $1-1/(2M)$. Variants bound decoders with $s$
modulation frequencies by rate $s/(s+1)$ and show that every address digit must be treated asymmetrically. The
hypothesis on the classes is needed: for $d=6$, two classes with no almost-period carry a code of dimension $25>24$
(\cref{prop:aperiodic}), found by a SAT solver and checked in Lean, and for every fixed number $M\ge2$ of classes the gain over
the bound is unbounded (\cref{sec:two-class}, added after a memo to this paper).
\end{itemize}

In words: to beat the pairing value, the decoder may have neither a hidden order nor a symmetry, on either index, that
the swap does not share. The second theorem uses exactly the feature of machines that cone relations forget.

\paragraph{What the theorems are not.} They are statements about the linear one-hop model. They are not lower bounds
for machines, and (T) remains open. \Cref{sec:construction} lists what they exclude and the three steps that would be
needed to carry them to machines; \cref{sec:small} gives exact values at small sizes; \cref{sec:open} states the
problems.

\paragraph{A subsequent memo, and the revisions of this paper.} After revision \texttt{b2c389a} of this paper in its
repository, a memo~\cite{ClassUniformMemo2026} proved, in Lean, that the two classes of \cref{prop:aperiodic} carry a
code of dimension $25$ over every field in which $3\ne0$; that class-uniform codes can be blown up, which gives gains
that grow with the number of classes; and bounds for the cost of advice in the one-hop model. Revision
\texttt{14b157a} of this paper cited the memo where it bears on the text, and corrected three passages:
\begin{itemize}[leftmargin=1.5em]
\item the charge for a bit that a genie reveals (\cref{sec:to-machines,prob:machines}): the earlier text charged every
revealed bit as a change of regime, at half a row, where advice that enters additively changes no decoder and half a
row is the price of a change that affects a whole row;
\item \cref{prob:aperiodic}(c), which the memo answered as it was worded and which was then stated for a fixed number
of classes;
\item the first sentence of \cref{sec:would-need}, which had dropped a hypothesis of the result it summarizes.
\end{itemize}
It added \cref{prop:entry,cor:entry-advice}, which state the corrected charge, with formal proofs of their bounds in
terms of the number of changes, and it updated \cref{rem:small}(d) and \cref{prob:aperiodic}(b).

An update of the memo then answered that form of \cref{prob:aperiodic}(c) for two classes: for every $d=4k\ge8$ it
gives an explicit rule, over fields in which $2\ne0$, whose code exceeds $(d^2+2d)/2$ by $d/4-1$, so that the gain of
two classes is unbounded. Its proof is written in the memo, which also proves in Lean, for $d=8$, $12$ and~$16$, that
the two classes carry a code of that dimension. It leaves open characteristic~$2$ and three or more classes.
Revision \texttt{ac70dc4} of this paper cited that update and added \cref{sec:two-class}, which settles both.
\begin{itemize}[leftmargin=1.5em]
\item \Cref{prop:two-class}: the memo's rule works over every field, $\F_2$ included, once its coefficients $2$ and
$\tfrac12$ are replaced by a nonzero $\lambda$ and its inverse.
\item \Cref{lem:mixed-blowup,cor:every-M}: a blow-up with a code of its own at each position turns the two-class codes
into codes for any fixed number $M\ge2$ of classes, with a gain of order $Md/8$ over $(d^2+Md)/2$. So the gain is
unbounded for every fixed number $M\ge2$ of classes, over every field.
\end{itemize}
These three results are proved in the text and are not formalized. \Cref{prob:aperiodic}(c) now asks for the rate of
growth.

Two papers of the series followed revision \texttt{ac70dc4}. One is about the original runs of adaptive
machines~\cite{AdaptiveTransfer2026}. The other is about the class-uniform rules of \cref{sec:symmetry}, their
codes and their composition~\cite{SharedDecoders2026}. Revision \texttt{3911d87} cites both where they bear on the text,
and changes the following.
\begin{itemize}[leftmargin=1.5em]
\item \emph{Bounds for classes of any shape} (\cref{prob:aperiodic}(a), \cref{sec:would-need}). For two classes the
second of these papers proves $\tfrac{41}{64}d^2+O(d^{3/2})$, where this paper had $\tfrac34d^2+O(d)$, and for $M$
classes a deficit of order $M^{-1/2}$, where this paper had $1/(2M)$. For partitions with a period, for which
\cref{thm:TP} shows that the gain is at most linear in~$d$, it determines the largest code of a lift under
hypotheses that are stated with \cref{prob:aperiodic}(c).
\item \emph{Which linearity the bounds need} (\cref{sec:sets}, new). Theorem~TP needs the code to be a subspace, but
not the decoders to be linear. For sets of matrices that are not subspaces it fails completely: one decoder for all
rows is correct on a fraction $1/(2d)$ of all matrices. Both facts are from that paper. That revision adds
\cref{prop:sets}: the elimination step of Theorem~FL and the bound for rules that copy one entry hold for sets and
arbitrary functions, as instances of the acyclic-subgraph bound of Bar-Yossef, Birk, Jayram and
Kol~\cite{BBJK2011}; and a remark on marker rows, by which a set of matrices with one function for all rows exceeds
the value $(d^2+d)/2$ of subspaces. One sentence of \cref{sec:symmetry-says} and one of \cref{sec:reading}, which said that
Theorem~TP needs linear decoders, are corrected, and a sentence is added to \cref{sec:symmetry-says}.
\item \emph{The step to machines} (\cref{sec:to-machines}, \cref{prob:machines}). The earlier text took the form of
its question from a count for affine rules. For decoders that are not affine and see the whole matrix there is no such
count. What machines add is that the cones a genie leaves are of bounded size, and whether shared decoders of
bounded window bound a set of matrices is the new \cref{prob:machines}(b). \Cref{sec:to-machines} also cites the
first bound for two hops, and notes that a perfect shuffle of the columns has $d-1$ classes; \cref{sec:fail} notes
what \cref{prop:sets} gives for synchronous access.
\item \emph{Adaptive machines} (\cref{sec:adaptive}, \cref{sec:one-function}, \cref{sec:exhaustions}). The first of
the two papers proves on the machine's own runs a statement of the form that \cref{sec:adaptive} reported as open,
and gives the function of \cref{sec:one-function} a form on those runs.
\item \Cref{prob:unequal} is now said to be about linear decoders.
\end{itemize}
\Cref{prop:sets} is proved in the text and is not formalized. The other results of this paper and their formal proofs
are unchanged, and so are their numbers.

That revision, \texttt{3911d87}, asked \cref{prob:machines}(b) for every window. The revision after it,
\texttt{a445cd8}, answered part of it, in two propositions added to \cref{sec:sets}.
\begin{itemize}[leftmargin=1.5em]
\item \Cref{prop:order}: the acyclic-subgraph bound, applied through an order of the rows or of the columns of the
matrix. For window two it gives at most $q^{d(d+1)/2+d-1}$ matrices, for all maps, all functions and all sets,
which is attained at $d=3$ and is within $d/2$ of an example at every even~$d$; and
for every window it gives $q^{(1/2+o(1))d^2}$ when the offsets, or the differences of the rows from a permutation,
take boundedly many values. Both hold for functions of whole rows, respectively whole columns.
\item \Cref{prop:mutual}: for every $\eps>0$, all but $\eps d^2$ of the pairs of positions that help each other are
related by one of boundedly many translates of the swap, so a code that uses only such pairs stays at the pairing
rate; and $o(d^2)$ positions meet all transitive triangles, which makes the theorems of the seventh paper for
relations without them available on one-class relations.
\item \Cref{prob:machines}(b) is now asked for windows of at least three, and \cref{sec:sets} reports computations
which indicate that for large windows the acyclic-subgraph bound cannot answer it.
\end{itemize}
The two propositions are proved in the text and are not formalized. No earlier result changes, and no number of a
result or of a table.

A further revision adds four propositions after them, on operators whose rows are arbitrary functions of whole rows
of their argument, and \cref{prob:machines}(c).
\begin{itemize}[leftmargin=1.5em]
\item \Cref{prop:rows}: such an operator transposes at most $q^{\sigma}$ matrices, where $\sigma$ is the least sum,
over the orders of the rows, of the place of the last row that each row of the output reads; this is exactly the
acyclic-subgraph bound of the relation, it equals $(d^2+sd)/2$ when each row of the output reads the block of $s$
rows that contains its image under a bijection, and it is at most $wd(d+1)/(w+1)$ when $w$ rows are read.
\item \Cref{prop:stages}: in a pipeline of stages that read rows inside cosets of the address space, with arbitrary
functions and alphabets between the stages, the dimensions of the cosets add up, and $t$ moved digits give at most
$q^{(d^2+2^td)/2}$ transposed matrices, which partial transposes attain. So a set of rate $\tfrac12+\delta$ needs
$\log_2d-\log_2(1/(2\delta))$ moved digits, for every code. Scans with a buffer of one row enter through advice,
and pipelines that mix differently arranged copies of the array through the number of ones of the address, with
$\sqrt{\log d}$ in place of $\log d$.
\item \Cref{prop:expanding}: if the rows read expand, the bound of \cref{prop:rows} is at least
$(\tfrac12+c)d^2$, and by the monotone expanders of Bourgain and Yehudayoff, in the form with three matchings of
Dujmovi\'c, Sidiropoulos and Wood, there are such relations made of the identity and three increasing maps, which
three forward passes gather when the strings that mark the maps are given. So the acyclic-subgraph bound cannot show
that operators of boundedly many such passes transpose only $q^{(1/2+o(1))d^2}$ matrices.
\item \Cref{prop:pieces}: products of codes on boundedly many positions stay at the pairing rate on every such
relation of bounded degree.
\end{itemize}
These are proved in the text and are not formalized, except that \cref{prop:expanding}(b) applies a cited theorem,
of which we use the statement and have not checked the proof. No earlier result changes, and no number of a result
or of a table.

\subsection{How to read this paper and its formal proofs}

The synthesis and the diagnosis contain no numbered results; they cite. The numbered results of
\cref{sec:order,sec:symmetry} each carry a status line.
\begin{itemize}[leftmargin=1.5em]
\item ``Lean:'' names the formal counterpart in the library \leanfile{lean/SwapRigidity/}; the Lean statement is
authoritative, and the text gives an outline.
\item ``Lean (partial):'' names what is formalized and says what is not.
\item ``Not formalized'' is followed by a complete proof.
\end{itemize}
Formalized in Lean~4 with Mathlib, on the standard axioms only: the coordinate form of Theorem~FL, its attainment and
the corollary for a pair of flags; the pigeonhole lemma, its composition property and the lower bound
$\dim\ker[\sw,T]\ge d$; Theorem~TP on a cycle with an arbitrary shift set, its attainment for equal pieces, and its
line form; the bound for decoders that change at single entries (\cref{prop:entry}), apart from a count, and its form
with advice (\cref{cor:entry-advice}); and the twist-class theorem. Not formalized, and proved in the text: the passage from flags to adapted bases; the corollaries
of Theorem~FL that need an algebraically closed field, the structure of $p$-groups or an arithmetic estimate
(\cref{cor:two-terms,cor:commutative,cor:affine,cor:expansion}); the equality case of \cref{prop:swap-floor}; the
copy-rule bound; the count in \cref{prop:entry}; \cref{cor:weyl}; \cref{cor:digit} apart from the abstract lemma it
rests on; the gains for a fixed number of classes over every field
(\cref{prop:two-class,lem:mixed-blowup,cor:every-M}); the bounds for sets
(\cref{prop:sets,prop:order,prop:mutual}) and for whole rows and stages
(\cref{prop:rows,prop:stages,prop:pieces} and \cref{prop:expanding}(a)); and the translation between class-uniform rules and operators $\sum_aA_a\otimes B_a$. Numerical observations appear only in remarks
labelled as numerical and in \cref{sec:small}, and no proof uses them. \Cref{app:leanindex} indexes the declarations and
\cref{app:verification} explains how to check them.

\subsection{Related work}

\paragraph{The series.} \Cref{sec:series} is the account of the earlier papers. The closest of them to the new theorems
are the eighth, which states the coding problem for linear machines as relative rigidity of the transposition
matrix~\cite{EighthPaper2026}; the second and fourth, which exclude locally affine head
motion~\cite{SecondPaper2026,FourthPaper2026}; and the seventh, whose elimination lemmas and Cayley corollary are the
storage-code relatives of \cref{thm:FL,thm:TP}~\cite{SeventhPaper2026}.

\paragraph{Matrix equations with a transpose.} Equations in which $X$ and $X^{\tr}$ both occur have been studied
as $\star$-Sylvester and $\top$-Stein equations. The solution space of $AX+X^{\star}B=0$ is described through the
Kronecker canonical form of a pencil~\cite{DeTeranDopicoGuilleryMontealegreReyes2013}; that of $XA+AX^{\tr}=0$ is
found in~\cite{DeTeranDopico2011} and applied to congruence orbits; unique solvability of $AXB+CX^{\star}D=E$ is
characterized in~\cite{DeTeranIannazzoPoloniRobol2018}, and of $X=AX^{\tr}B+C$ in~\cite{Chiang2013}. Generalized Schur
decompositions are used for solvability and numerical solution~\cite{ChiangChuLin2012}. \Cref{thm:FL} uses
triangular forms for a different purpose, a dimension bound in the singular case with one side unconstrained. We found no statement of \cref{cor:two-terms}, of the
bound for an arbitrary number of terms with one triangularizable side, or of any bound on the rank of the commutation
matrix~\cite{MagnusNeudecker1979} minus a structured sum of Kronecker products. Our search was limited to abstracts and
indexes, and we claim no priority for statements that specialists may regard as folklore.

\paragraph{Commutators and network codes.} \Cref{lem:pigeonhole} is the fact that a commutator has rank at most twice
the rank of either factor; we know no source that states it in the form used here and we give its proof. Transposition
as a multi-commodity network-coding instance goes back to~\cite{AHJKL2006}. Network codes whose local kernels are sums
of cyclic shifts have been studied for multicast~\cite{TangSunLiYangLong2019}; we found nothing on shift-covariant
decoders for transposition.

\subsection{Organization}

\Cref{sec:prelim} fixes the problem and the two models. \Cref{sec:series} is the synthesis and \cref{sec:wall} the
diagnosis. \Cref{sec:order,sec:symmetry} prove the two theorems. \Cref{sec:construction} is the construction side,
\cref{sec:small} the exact computations and \cref{sec:open} the open problems.

\section{The problem and two models}\label{sec:prelim}

\subsection{The question}\label{sec:question}

Throughout, $d\ge1$ is the side length of a square matrix, $b(d)=\max(1,\clog d)$ is the entry width, and $L=d^2b(d)$ is
the input length. A machine is a deterministic multitape Turing machine with a read-only input tape, a fixed number~$k$ of
work tapes, a write-only output tape and a fixed finite alphabet. The input is a $d\times d$ matrix with entries in
$\{0,1\}^{b(d)}$, written row by row without delimiters; the output is its transpose in the same encoding.

\begin{quote}
\textbf{(T)} There are a machine $M$ and a constant $c$ such that, for every $d$ and every input matrix, $M$ halts with
the transpose on its output tape within $c(L+1)$ steps.
\end{quote}

The authoritative statement is the Lean definition \texttt{Transpose.LinearTimeTransposable} of the first
paper~\cite{FirstPaper2026}, which the series keeps frozen. A machine is \emph{oblivious} if its head positions at each
step, the steps at which it writes output and its halting time depend only on the input length; otherwise it is
\emph{adaptive}. A \emph{pass} is a maximal stretch of steps on
which every head moves monotonically.

The folklore algorithm transposes in $O(L\log d)$ steps by a recursion on address digits: write the row address and the
column address of an entry in binary, and exchange them one digit at a time, each exchange costing $O(L)$. An unrefereed
preprint~\cite{OpenAI2026mult} claims $O(L(\log d)^{1-\kappa})$ for a small $\kappa>0$. Harvey and van der
Hoeven~\cite{HvdH2025} reduce transposition to integer multiplication on multitape machines, so a linear-time multiplier
would give~(T), and a proof that (T) fails would be a superlinear lower bound for off-line multiplication on these
machines. No superlinear lower bound is known for transposition, nor for off-line multiplication, sorting or the
discrete Fourier transform, on machines with two or more work tapes. The known bounds are for one work
tape~\cite{DMS1991,DietzfelbingerMaass1993}, for machines that only move data~\cite{Stoss1973}, and for
reversal-bounded conservative machines~\cite{Kirchherr1988}; the superlinear bounds that follow from the separation of
deterministic and nondeterministic linear time~\cite{PPST1983} concern problems that are hard for nondeterministic
linear time~\cite{Grandjean1988}.

\subsection{Storage codes on recovery relations}\label{sec:storage}

Most of the series is written in the following language, which we use in \cref{sec:series,sec:wall}.

\begin{definition}[Storage codes]\label{def:storage}
A \emph{recovery relation} on a finite set $V$ assigns to every $v\in V$ a set $N(v)\subseteq V$; it is
\emph{loopless} if $v\notin N(v)$ for all~$v$. A
\emph{storage code} on it, over an alphabet $\alpha$, is a set $C\subseteq\alpha^V$ such that for every $v$ the coordinate
$w_v$ of a codeword $w\in C$ is a function of $(w_u)_{u\in N(v)}$. Its \emph{rate} is $\log\card C/(\card V\log\card\alpha)$.
\end{definition}

A relation in which $u\in N(v)$ exactly when $v\in N(u)$ is a graph, and the largest $\log_{\card\alpha}\card C$ is then the
guessing number of Riis~\cite{Riis2007}; storage codes are dual to index codes with side information~$N$. On a perfect
matching the rate is exactly $\tfrac12$: each vertex copies its partner. We call $\tfrac12$ the \emph{pairing rate}.

The \emph{incidence graph} of a relation has two vertices $L_p$ and $R_p$ for every coordinate~$p$, its port as a
helper and its port as a receiver, and an edge $L_pR_q$ whenever $p\in N(q)$. The \emph{paired distance} of $p$ is the
distance between $L_p$ and $R_p$ in this graph; for a graph it is the length of the shortest odd closed walk through~$p$.
Following the fifth paper~\cite[Definition~2.3]{DirectedRecovery2026}, we call a sequence of relations \emph{separated}
if, for every fixed~$r$, all but $o(\card V)$ coordinates have paired distance greater than~$r$. A perfect matching has
all paired distances infinite, so separation does not exclude the pairing; it excludes gains that come from short
cycles.

For an oblivious machine and a set of steps whose values are revealed (a \emph{genie}), the \emph{cone} of an output
bit is the set of input bits on which it can still depend through the time--space dependency graph. The \emph{cone
relation} assigns to an input coordinate $q$ the cone of the output bit that must equal it. If the machine transposes,
then for some value of the revealed word the inputs that agree with it form a storage code on the cone relation, of
high rate when the genie reveals few bits (\cref{sec:depth}). The \emph{depth-$D$ ball} of an output bit is the set of
input positions read at steps within $D$ dependency steps of the step that writes it, with no genie.

\subsection{The linear one-hop model}\label{sec:onehop}

The two theorems of this paper live in a smaller model, in which the code and the geometry are visible at once. Let $F$
be a field, $V=F^d$, and let $\sw$ be the swap $u\otimes w\mapsto w\otimes u$ on $V\otimes V$. Under the identification
of $V\otimes V$ with the $d\times d$ matrices $\Mat_d(F)$, $\sw$ is transposition, $X\mapsto X^{\tr}$, and an operator
\[
  K=\sum_{a}A_a\otimes B_a,\qquad A_a,B_a\in\Mat_d(F),
\]
acts by $X\mapsto\sum_aA_aXB_a^{\tr}$. Every operator on $V\otimes V$ has this form for some finite family of terms.

\begin{definition}[Swap code]\label{def:swapcode}
\formal{Flag.lean}{swapCode,mem_swapCode}
The \emph{swap code} of $K=\sum_aA_a\otimes B_a$ is
\[
  \ker(\sw-K)=\Bigl\{X\in\Mat_d(F): X^{\tr}=\sum_aA_aXB_a^{\tr}\Bigr\}.
\]
Its \emph{rate} is $\dim\ker(\sw-K)/d^2$. (The Lean definition is the set of matrices on the right; its identification
with the kernel of an operator is not used in the formal proofs.)
\end{definition}

The code is the set of matrices that $K$ transposes correctly. Read row by row, the equation says that every column
of~$X$ is a fixed linear function of the rows of~$X$: the stored word is $X$ in row-major order, the receivers want it in
column-major order, and $K$ is the decoder. For $K=1\otimes1$ the code is the space of symmetric matrices, of dimension
$d(d+1)/2$. We call this the \emph{pairing value}; its rate tends to the pairing rate~$\tfrac12$. For $K=\sw$ the code
is everything.

Entries of $b$ bits are handled by replacing $V\otimes V$ with $V\otimes V\otimes F^b$ and $\sw$ with $\sw\otimes1$; the
flag bounds below then hold with an extra factor~$b$ (\cref{rem:entries}), and \cref{thm:TP} likewise.

\paragraph{Relative rigidity.} A bound $\dim\ker(\sw-K)\le(1-\rho)d^2$ for all $K$ in a class says that the transposition
matrix is $\rho d^2$-rigid relative to that class: $\sw$ is not the sum of a member of the class and a matrix of rank
below $\rho d^2$. The eighth paper shows that, for machines with linear data paths, the coding half of the series'
programme is exactly such a statement, for the bit-level transposition matrix and the class of matrices whose row
for each output bit is supported on its depth-$D$ ball in a
schedule~\cite[Theorem~7.5 and Problem~10.6]{EighthPaper2026}; this is a relative form of Valiant's
rigidity~\cite{Valiant1977}. The classes of this paper are defined by algebraic structure on one tensor side instead of
by a sparsity pattern. The other side is arbitrary.

\section{The series in one place}\label{sec:series}

This section states the main results of the sixteen earlier papers and their eight memos, grouped by what they are about
and not by when they appeared. Nothing is reproved. Each statement carries its citation and its main hypotheses; the
cited theorems have further parameters (the number of tapes, the alphabet, the time constant), which we suppress.
\Cref{tab:series} at the end lists the papers in order. A reader who wants only the conclusion can skip to
\cref{sec:wall}.

\subsection{The question, passes and the coding formulation}\label{sec:first}

The first paper~\cite{FirstPaper2026} fixes the formal statement of~(T), shows that for linear time it agrees with the
simplest case of the formulation of Harvey and van der Hoeven~\cite[Proposition~11.1, a written proof]{FirstPaper2026},
and proves the following.
\begin{itemize}[leftmargin=1.5em]
\item \emph{Passes.} A transposer with at least one work tape makes $\Omega(\log d)$ passes, whatever it does with its
data: for oblivious machines with no premise on the running time, and for all machines that run in time polynomial
in~$L$~\cite[Theorems~5.9 and~5.12]{FirstPaper2026}. A memo removes the runtime premise for halting machines that are
correct on all inputs~\cite{PassesMemo2026}. This is the only logarithmic lower bound on a resource of the machine itself
that the series proves for all machines. It counts passes, not steps, and it is tight: the folklore algorithm uses
$O(\log d)$ passes.
\item \emph{Folding.} Along its popular traces, every limit of the rescaled work-head motion of a linear-time
transposer satisfies an isoperimetric inequality that forces the heads to fold back
(Theorem~M*,~\cite[Theorems~7.8 and~7.9]{FirstPaper2026}). The proof is a cut argument: the information that crosses a
cut of the time--space diagram is bounded by the states and the work-tape cells that enter across the cut.
\item \emph{Coding.} Every linear-time transposer yields a causal network code for complete exchange on the block
network of a time--space diagram of linear length, correct on a $2^{-O(d\log d)}$ fraction of the inputs, and on all
inputs if the machine is oblivious~\cite[Theorems~8.10 and~9.9]{FirstPaper2026}. The \emph{Tape Coding Conjecture} says
that for large $d$ no causal code solves complete exchange on such a fraction of the inputs on the block network of any
diagram of linear length; it implies that (T) fails~\cite[Conjecture~9.1, Theorem~9.2]{FirstPaper2026}. Its plain form,
for codes correct on all inputs, excludes oblivious linear-time transposers~\cite[Theorem~9.9]{FirstPaper2026}. On
these networks every routing has rate $O(1/\log d)$, so a solving code beats routing by a factor $\Omega(\log d)$, and
the conjecture is equivalent to saying that coding beats routing by only $o(\log d)$ on tape
networks~\cite[Theorems~8.8 and~9.4]{FirstPaper2026}.
\end{itemize}
The pass bound at polynomial running time, Theorem~M* and the reduction to codes survive when the machine is correct
on only a $2^{-O(d\log d)}$ fraction of the inputs. The same paper studies an abstract model of the one technique claimed
to beat $O(L\log d)$, recursion on XOR-mixed address frames: there every gadget costs at least its number of wires,
whereas linear time would need less~\cite[Theorem~12.4 and Proposition~12.9]{FirstPaper2026}; the correspondence between
the model and the preprint's recursion is informal. It also records why a proof must use the shape of tapes: a bound
$o(\log n)$ on coding gains for bounded-degree graphs in general would give superlinear circuit lower bounds for cyclic
shifts and multiplication~\cite[Proposition~11.2]{FirstPaper2026}.

The third paper~\cite{ThirdPaper2026} shows that the bit case is at least as hard: linear-time transposition of bit
squares of dyadic side implies~(T), by one fixed wrapper. The converse is open, and the paper proves barriers for
restricted compilers in that direction; one of them is a multiscale obstruction for compilers with boundedly many
stages.

\subsection{What head motion must do}\label{sec:motion}

A second line asks what the heads of a linear-time transposer must do.
\begin{itemize}[leftmargin=1.5em]
\item \emph{Not affine.} Along the popular traces, the limits of the rescaled work-head paths of a linear-time
transposer fail to be locally affine on sets of times of total measure at least $\tfrac12$, summed over the heads and
measured in blocks. Bounded-length words over three explicit palettes of scan operations are excluded, and so are
contracting recursions all of whose stages have regular limiting
motion~\cite[Theorems~6.3, 7.1 and~7.5]{SecondPaper2026}. No machine whose work heads are home-based with legs made of
boundedly many constant-rate runs transposes in linear time, and in every rescaled limit a linear-time transposer has
a linear amount of curved travel~\cite[Theorems~3.10 and~3.12]{FourthPaper2026}.
\item \emph{Macroscopic crossings.} An oblivious linear-time transposer routes a linear amount of information through
pairs of interleaved revisits whose excursions are a fixed fraction of~$L$: for every set of at least $\alpha L$ output
bits, every set of steps separating them from the earlier ends of such crossings has at least $c$ times as many
elements~\cite[Theorem~5.5]{SixteenthPaper2026}. A fixed fraction of this expansion sits on one coarsely monotone leg of
one tape, in a skeleton from a finite family independent of~$d$, and two work heads must each make linearly many long
revisits~\cite[Theorems~6.7 and~4.1]{SixteenthPaper2026}. The crossing and skeleton statements hold on heavy traces of
adaptive machines and of a bit simulator~\cite[Theorem~8.1]{SixteenthPaper2026}.
\item \emph{Rows to columns.} In the diagram of an oblivious transposer, every set of steps that separates the inputs
of a set of rows from the outputs they feed into a set of output rows carries the corresponding information, whatever
the coding~\cite[Lemma~7.1]{SixteenthPaper2026}. When the cuts for different address bits are disjoint this gives a
time bound $\Omega(L\log d)$ that no coding can evade~\cite[Corollary~7.3]{SixteenthPaper2026}.
\item \emph{Source covers.} For one fixed geometry of returns, a certificate that retains every input on which an
output can depend, and charges the information used to simplify the dependence, gives a correctness deficit of at
least $\eps L$ when the return families consist of a two-class core, at most one arbitrary exceptional family and
families of at most $\eps L$ hops, for fixed bounds on the running time, the capacity and the number of span
classes~\cite[Corollary~4.3]{SourceCovers2026}. A memo excludes funnels and single-sweep terminal landings and shows
where the certificate stops~\cite{GapAreaMemo2026}. These hypotheses on the geometry are not derived for arbitrary
machines.
\end{itemize}

Three results mark the limits of this line.
\begin{itemize}[leftmargin=1.5em]
\item \emph{Padding.} For every machine there is a wrapper, with more tapes and states and additive overhead
$O(d^2b)$, whose complete runs satisfy the cut inequality behind Theorem~M* with zero slack, and whose normalized head
limits satisfy all the conclusions of Theorem~M* under a linear time bound. The wrapper's motion is independent of the
computation that follows~\cite[Theorem~8.1]{SecondPaper2026}. It needs no idle tapes, and it also meets the finer
entry-cell form of the inequality at every slot length from one cell up to a constant number of
rows~\cite{PaddingMemo2026}. So these conditions on motion can be met by a machine that transposes nothing; they do not
see which information reaches the outputs.
\item \emph{One cut for all pairs.} In the diagram of any run, one set of steps separates the inputs from the outputs
for every pair of row sets at once~\cite[Proposition~10.2]{SixteenthPaper2026}. So counting the cuts of the
rows-to-columns inequality cannot by itself go beyond linear time: the disjointness that
\cite[Corollary~7.3]{SixteenthPaper2026} needs is not automatic.
\item \emph{The candidates.} Five explicit linear-time cascade schedules meet the crossing and skeleton conditions
at the sizes tested. All five are excluded, for fixed work alphabets, by the regular-recursion obstruction of the second
paper~\cite{CascadeObstructionMemo2026}, and two of them also by explicit cuts of the rows-to-columns
inequality~\cite[Theorem~10.8]{SixteenthPaper2026}. Whether the combined conditions on motion decide the oblivious case
is open.
\end{itemize}

\subsection{Depth reduction}\label{sec:depth}

The fourth paper~\cite{FourthPaper2026} splits the oblivious case of~(T) into two statements.
\begin{itemize}[leftmargin=1.5em]
\item \emph{Depth reduction} (the trajectory cut problem $\TCP(k)$): for every $\eps>0$, deleting $\eps T$ edges from
the dependency graph of $k$ head walks of length~$T$ leaves no path with more than a bounded number of long hops, and
then, after periodic steps are also revealed, paths of bounded length.
\item \emph{The crux}: an index-coding lower bound on the resulting bounded cones, strong enough to exceed the cost of
the revealed steps.
\end{itemize}
Together they would exclude oblivious linear-time
transposers~\cite[Theorem~4.10 and Corollary~4.11]{FourthPaper2026}. Depth reduction is now a theorem. It was proved
for home-based walks~\cite[Theorem~4.20]{FourthPaper2026}; then for walks of bounded nesting load, and for one arbitrary
head together with heads of bounded nesting load at boundedly many
scales~\cite[Theorems~3.6 and~4.2]{EighthPaper2026}; and finally for every number of walks, with a depth $D(k,\eps)$
such that $\log\log D=O(k^2\log(k/\eps))$~\cite[Theorem~7.5]{BoundedStack2026}. A memo shows that $\TCP(k)$ for all~$k$
is equivalent to constant-depth revelation of arbitrarily small entropy per input bit, for all linear-time oblivious
computations~\cite{InfoDepthMemo2026}. The fourteenth paper also refutes Conjectures~A$'$ and~A$''$ of the fourth and
eighth papers, which would have made the deep-nesting pairs
depth-robust~\cite[Corollary~9.1]{BoundedStack2026}: walks nested at $m$ scales are cut to a depth independent of~$m$.
It draws the conclusion that matters here~\cite[Corollary~9.3]{BoundedStack2026}:
\begin{quote}
\emph{for oblivious machines, the crux alone would exclude linear-time transposers.}
\end{quote}
In the language of \cref{sec:storage}: for an oblivious linear-time transposer and every $\eps>0$, revealing $\eps L$
bits at suitable steps leaves every output bit a function of the revealed word and of boundedly many input bits, and
for some value of the revealed word the inputs compatible with it form a binary storage code of rate at least $1-\eps$
on the machine's own cone relation. The same genies give depth reduction at the level of
codes~\cite[Corollary~9.6]{BoundedStack2026}, as conjectured in the eleventh paper. And, given the \emph{native half
bound} (every binary storage code on the cone relations of the machine's own traces has rate at most
$\tfrac12+o(1)$, which is open), they exclude every linear-time transposer, adaptive or not, whose trace has entropy at
most $\eta L$ for a fixed $\eta<\tfrac12$~\cite[Corollary~9.7]{BoundedStack2026}.

For machines whose data path is linear over a field, linear depth reduction together with a \emph{linear crux}, a
min-rank bound for the ball systems of schedules, excludes oblivious linear-time
transposers~\cite[Theorem~7.5]{EighthPaper2026}. The linear crux is exactly the relative rigidity of the transposition
matrix described in \cref{sec:onehop}~\cite[Problem~10.6]{EighthPaper2026}, and since the fourteenth paper it alone
suffices~\cite[Corollary~9.3(c)]{BoundedStack2026}.

\subsection{Storage codes: what is proved}\label{sec:proved}

The crux needs an upper bound on the rate of storage codes on cone relations. The series proved such bounds, at the
pairing rate, for many classes of relations.
\begin{itemize}[leftmargin=1.5em]
\item \emph{Homomorphisms and small degree.} Rate bounds pull back along homomorphisms, so every storage code on a
graph that maps to $C_5$, to $K_{14/5}$ or to the Petersen graph has rate at most
$\tfrac12$~\cite[Theorem~6.5 and Corollary~6.9]{FourthPaper2026}. On triangle-free graphs of maximum degree three no
scalar linear code over any field beats~$\tfrac12$~\cite[Theorem~6.13]{FourthPaper2026}. The Clebsch graph, of degree
five, has guessing number~$10$ on $16$ vertices~\cite[Corollary~6.11]{FourthPaper2026}, so bounded degree alone is not
enough.
\item \emph{Bounded complex coefficients.} On locally bipartite graphs of bounded degree, linear codes with bounded
complex coefficients have rate at most $\tfrac12+o(1)$~\cite[Theorem~6.17]{FourthPaper2026}. The same holds for complex
linear recovery with bounded coefficients and fixed block width on directed relations of bounded incidence degree that
are separated, hence on cone relations after a paid trimming of
degrees~\cite[Theorem~4.5 and Corollary~4.7]{DirectedRecovery2026}. On such relations a gain by linear codes with
boundedly many symbols per vertex therefore needs positive characteristic, or complex coefficients that no diagonal
rescaling makes bounded.
\item \emph{Hyperfinite relations.} Every code over every finite alphabet on relations whose incidence graphs are
uniformly hyperfinite and separated has rate at most $\tfrac12+o(1)$; no degree bound is
needed~\cite[Corollary~4.4]{DirectedRecovery2026}.
\item \emph{Nested elimination.} The exact bound $\card C^2\le\card\alpha^{\card V}$ holds for all codes, linear or not,
on large classes of relations without transitive triangles: triangle-free graphs of maximum degree three; loopless
relations whose underlying graph has maximum degree three; relations whose vertices of in-degree at least three induce
an acyclic sub-relation; relations partitioned into blocks with private outside parents, each block chordless
subcubic, a chorded digon path or internally certified; and Cayley relations $\mathrm{Cay}(G,\{s,s^{-1},t\})$ of every
finite group~\cite[Theorems~4.3, 5.2, 5.3 and~5.10, Corollary~5.12]{SeventhPaper2026}. It also holds on the class of
the fourth paper's Theorem~S, for all codes~\cite[Theorem~3.6]{DirectedReciprocity2026}. The degree bound cannot be
raised: a $4$-regular triangle-free graph on $13$ vertices carries a linear code over $\F_4$ with
$\card C^2=4^{14}$~\cite[Example~7.8]{SeventhPaper2026}.
\item \emph{Covers and towers.} For the binary adjacency code of the Clebsch graph, every connected tower of binary
lifts has kernel dimension at most $N/2+2$~\cite[Theorem~9.2]{Covers2026}; lifting a binary linear code with a
unit-diagonal check along a tower of double covers never raises its rate~\cite[Theorem~8.1]{EighthPaper2026}.
\item \emph{Memos.} The half bound for circular cliques $K_{p/q}$ extends to every $p/q<3$~\cite{CircularMemo2026};
nested elimination extends to soft recovery with record capacities~\cite{SoftRecovery2026}; and the bounds for supports
of maximum degree three extend beyond scalar linear codes, to reversible rules over any alphabet and to linear codes
whose rows have at most one port of rank above one~\cite{SubcubicMemo2026}.
\end{itemize}

One result of this group is about the relations of machines themselves. For cones of fixed depth in linear-time
schedules, after a paid trimming of degrees, only $O(L/\log d)$ matched pairs of an input port and the output that
demands it lie within any fixed incidence
distance~\cite[Theorem~3.3, Corollaries~3.7 and~3.8]{DirectedRecovery2026}. So cone relations are separated in the
sense of \cref{sec:storage}, and almost no output is within bounded distance of its own input bit. No gadget of bounded
size can then lift a code above the pairing rate: a gain above $\tfrac12$ on boundedly many coordinates needs a short
path between the two ports of one of them~\cite[Lemma~4.1 and Theorem~4.2]{DirectedRecovery2026}. The $L/\log d$ comes
from a bound in steps for transport at bounded occupancy, which charges every transported label the $\log d$ bits of
its row address~\cite[Lemma~3.4]{DirectedRecovery2026}: for uncoded movement the logarithm is proved, and what the
bound does not cover is a cell that carries a combination of many labels.

This led to the \emph{all-code question}~\cite[Question~9.1]{DirectedRecovery2026}: does every storage code on
separated relations of bounded incidence degree have rate at most $\tfrac12+o(1)$?

\subsection{Storage codes: what is false}\label{sec:false}

The answer is no.
\begin{itemize}[leftmargin=1.5em]
\item \emph{Above the pairing rate.} There are relations of bounded (astronomically large) incidence degree in which
every paired distance tends to infinity, with $\F_{127^2}$-linear storage codes of rate above $0.5021$, and binary
codes of rate above $0.5013$; for large primes inert in $\Q(\sqrt5)$ the rates approach
$4-2\sqrt3\approx0.536$~\cite[Theorems~7.1 and~7.3, Corollary~8.2]{StorageCodes2026}. The relations are twisted
quotients of a Bruhat--Tits tree, the codes are Tanner codes with Reed--Solomon local codes, and their dimension comes
from Riemann--Roch on Hilbert modular surfaces in characteristic~$p$. The undirected form of the question, on graphs of
bounded degree with long odd cycles, has a negative answer as well~\cite[Corollary~9.3]{StorageCodes2026}. Tanner operators of the same
kind, at $p=3$, refute the independence of the chain in Jaikin-Zapirain's arbitrary-field form of the L\"uck
approximation conjecture~\cite[Theorem~5.6]{Luck2026}.
\item \emph{Rates approaching one.} A linear-time transposer, even an adaptive one, would yield binary storage codes
of this kind at every large degree cap~$\Delta$, with rates tending to one as $\Delta$
grows~\cite[Theorem~3.1]{HighRate2026}; and such codes exist abstractly, with rate at least
$1-C\Delta^{-\alpha}$~\cite[Theorem~5.5]{HighRate2026}. So, at the scale of that transfer theorem, no rate bound for
separated relations of bounded degree in general can exclude a transposer~\cite[Corollary~6.2]{HighRate2026}: a lower
bound must use a further restriction imposed by the physical computation.
\end{itemize}
No tape process is known that realizes the recoveries of either family.

\subsection{Adaptive machines}\label{sec:adaptive}

Much of the above already covers adaptive machines: the pass bound (at polynomial running time, or by the memo),
Theorem~M* along popular traces, and the reduction of~(T) to the Tape Coding Conjecture. Beyond that the series has the
following.
\begin{itemize}[leftmargin=1.5em]
\item Every linear-time transposer yields an external-memory transposition with $O(d)$ I/Os on blocks of one row that
is oblivious on $2^{L-O(d)}$ inputs, so an external-memory lower bound without the indivisibility assumption would
refute~(T)~\cite[Theorem~6.4 and Corollary~6.5]{MergeDepth2026}.
\item The source-cover exclusions extend to machines of sufficiently low
adaptivity~\cite[Theorem~6.1]{MergeDepth2026}.
\item Given the native half bound for the relations of skeleton classes, which is open, a linear-time transposer whose
work heads enter every cell boundedly often must carry about $L/2$ bits in its coarse skeleton, hence in the coarse
motion of its input head~\cite[Theorem~4.13 and Corollary~4.16]{MergeDepth2026}; the fourteenth paper removes the
hypothesis on the work heads from the skeleton statement~\cite[Corollary~9.7]{BoundedStack2026}.
\item The fifteenth paper extracts its codes from any linear-time transposer through a packet
compilation~\cite[Theorem~3.1]{HighRate2026} and, at the original bit positions, from the traces of a bit simulator
with few traces~\cite[Theorems~8.1 and~9.2]{HighRate2026}. The crossing theorems hold on heavy traces of the machine and
of that simulator~\cite[Theorem~8.1]{SixteenthPaper2026}; whether they hold in the machine's own diagram on traces that
carry almost all inputs was left open there~\cite[Problem~11.2]{SixteenthPaper2026}. A paper that followed this one
proves a statement of that form at the canonical entry width: on the original run of all but an exponentially small
share of the inputs that a linear-time transposer handles, every cut between the outputs of a large set and the
macroscopic crossings has size proportional to that set~\cite[Theorem~H.7]{AdaptiveTransfer2026}. For every entry
width, and again outside an exponentially small share of the inputs, it proves that two work heads make linearly many
macroscopic returns, and that cuts between a large set of outputs and the reads of their inputs are of linear
size~\cite[Theorem~4.3 and Corollary~5.4]{AdaptiveTransfer2026}.
\end{itemize}

\begin{table}[t]
\centering\small
\begin{tabular}{@{}r>{\raggedright\arraybackslash}p{0.27\textwidth}>{\raggedright\arraybackslash}p{0.62\textwidth}@{}}
\toprule
 & Paper & Main contribution to the question \\
\midrule
1 & Tape Coding Conjecture~\cite{FirstPaper2026} & Formal statement; $\Omega(\log d)$ passes; Theorem~M*; transposers
are network codes \\
2 & Local Coding Barriers~\cite{SecondPaper2026} & Local coding bounds; affine and regular motion excluded; padding \\
3 & Bit and Word~\cite{ThirdPaper2026} & Transposition of dyadic bit squares implies~(T) \\
4 & Curved Motion~\cite{FourthPaper2026} & Curved motion; depth reduction plus crux; storage codes at rate $\tfrac12$ \\
5 & Directed Recovery~\cite{DirectedRecovery2026} & Paired-distance separation of cones; half bounds; the all-code
question \\
6 & Covers~\cite{Covers2026} & Half bounds along covers and for locally generated codes \\
7 & Nested Elimination~\cite{SeventhPaper2026} & Exact all-code half bounds by elimination \\
8 & Bounded Nesting~\cite{EighthPaper2026} & Depth reduction at bounded nesting; the linear crux as rigidity \\
9 & Source Covers~\cite{SourceCovers2026} & Source-cover certificates and exclusions \\
10 & Directed Reciprocity~\cite{DirectedReciprocity2026} & All-code bound on the class of Theorem~S; deep nesting \\
11 & Merge Depth~\cite{MergeDepth2026} & Merge depth; adaptive machines; external memory \\
12 & L\"uck Approximation~\cite{Luck2026} & A Tanner operator whose kernel survives along a congruence tower and dies
along pro-$3$ towers \\
13 & Above Half Rate~\cite{StorageCodes2026} & The all-code question has a negative answer \\
14 & Bounded-Stack DAGs~\cite{BoundedStack2026} & Depth reduction for every number of walks \\
15 & High-Rate Codes~\cite{HighRate2026} & A transposer needs rates tending to one; such codes exist abstractly \\
16 & Macroscopic Crossings~\cite{SixteenthPaper2026} & Crossings, skeletons, rows to columns, one cut for all pairs \\
\bottomrule
\end{tabular}
\caption{The earlier papers of the series.}\label{tab:series}
\end{table}

\section{Where the wall is}\label{sec:wall}

This section is a diagnosis. It contains no new mathematics. Its claims are consequences of the results cited in
\cref{sec:series}, and where it gives a reading instead of a consequence it says so.

\subsection{What is left is one function}\label{sec:one-function}

For oblivious machines the series has reduced~(T) to a question about codes (\cref{sec:depth}). Fix the number of
tapes, the alphabet and the time constant, and let $\eps>0$. Depth reduction gives, for every linear-time schedule, a
genie that reveals at most $\eps L$ bits and leaves cones of depth at most $D_\eps$, a bound that depends on $\eps$ and
on the fixed parameters but not on~$d$. Write $r(\eps)$ for the best rate, in the limit of large~$d$, of a binary
storage code on the cone relations that such genies leave on such schedules. Then:
\begin{itemize}[leftmargin=1.5em]
\item an oblivious linear-time transposer needs $r(\eps)\ge1-\eps$ for every $\eps>0$;
\item a proof that $r(\eps)<1-\eps$ for one $\eps$, for each choice of the fixed parameters, would exclude oblivious
linear-time transposers~\cite[Corollary~9.3]{BoundedStack2026};
\item the pairing gives $r(\eps)\ge\tfrac12$, and whether $r(\eps)>\tfrac12$ for any~$\eps$ is open in both directions.
\end{itemize}
The depth bound matters: without it the cone of an output is everything read before it, and the question is empty.
The only general upper bound is the elementary one for relations of degree at most~$\Delta$, a rate of
$1-1/(\Delta+1)$~\cite[Proposition~4.14]{FourthPaper2026}, which is useless at the depths in question. For machines with linear data paths the same reduction
ends in the relative rigidity of the transposition matrix on ball patterns (\cref{sec:depth}). For adaptive machines
the question arises in two forms: the Tape Coding Conjecture of the first paper, which needs no simulator, and
the storage-code criteria on the traces of a bit simulator (\cref{sec:adaptive}).

For particular classes of schedules the exclusion is proved. The sixteenth paper proves it for schedules whose
large output sets can be cut cheaply from the macroscopic crossings, and for schedules with boundedly many macroscopic
phases~\cite[Theorems~5.4, 9.2 and~9.3]{SixteenthPaper2026}; the source covers of the ninth paper exclude restricted
geometries for all codes. What is missing is a statement for all schedules.

\subsection{Three kinds of argument and their limits}\label{sec:exhaustions}

\paragraph{Motion without data.} The cut inequality behind Theorem~M*, its entry-cell refinement and every conclusion
of Theorem~M* are satisfied by the padded machine of the second paper, which transposes nothing (\cref{sec:motion}).
So conditions on head motion that do not refer to the outputs are necessary and cannot be decisive. Conditions that
do refer to them go further: the rows-to-columns inequality of the sixteenth paper excludes two of the five candidate
schedules. But one cut serves every pair of row sets, so this inequality gives a logarithm only when the cuts for
different address bits are disjoint, which is a hypothesis and not a theorem; and what excludes all five candidates is
the regular-recursion obstruction, which applies only to schedules all of whose stages have regular limiting motion.
Whether conditions of this kind decide the oblivious case is open.

\paragraph{Codes without machines.} On abstract separated relations of bounded degree, storage codes exceed the pairing
rate, and their rates tend to one polynomially fast as the degree cap grows (\cref{sec:false}). A transposer is only
shown to need rates that tend to one very slowly, and the cap grows as $\eps$ falls because a genie of cost $\eps L$
has depth at least of order $1/\eps$~\cite[Proposition~4.13]{FourthPaper2026}. So no rate bound for separated relations of bounded
degree in general can separate a transposer from an abstract code. The two properties of cone relations that the
series' coding bounds use, bounded degree after trimming and paired-distance separation, are shared by the
counterexamples.

\paragraph{Counting passes.} The pass bound of the first paper is the one logarithm proved for a resource of every
machine. It saturates one scale at a time: a linear-time run has room for about $2^j$ legs at scale $2^{-j}L$, while
counting at a single scale asks for only $O(j)$ of them, and the threshold of the coarse-pass bound of the sixteenth
paper is tight~\cite[Theorem~10.3 and Corollary~10.4]{SixteenthPaper2026}. That paper states the problem of an
inequality in the numbers of legs at all scales at once~\cite[Problem~11.1]{SixteenthPaper2026}; nothing in the series
proves one.

\subsection{Why every general tool saturates}\label{sec:saturation}

The $\log d$ of the folklore algorithm is a sum over $\log_2d$ address digits. Exchanging one digit of the row address
with one digit of the column address costs $O(L)$, and the transposition is the product of these $\log_2d$ commuting
exchanges. Linear time pays for any one of them in full, and for any bounded number.

Every tool of the series \emph{that applies to all machines} is single-scale or scale-free. Cut arguments, macroscopic
limits and coarse passes look at one scale at a time, and a linear-time machine has enough steps to satisfy any
single-scale requirement. Storage-code bounds are scale-free: they see a relation, not the scale of the addresses that
produced it. The one attempt to make a geometric quantity grow with the number of scales failed: walks nested at $m$
scales are cut to a depth independent of~$m$ (\cref{sec:depth}).

The series does charge scales separately in restricted settings, and they show what the general case lacks.
\begin{itemize}[leftmargin=1.5em]
\item For transport at bounded occupancy, each moved label is charged its $\log d$ address bits, in
steps~\cite[Lemma~3.4]{DirectedRecovery2026}. What the bound does not cover is coding.
\item For traces whose cuts for different address bits are disjoint, the rows-to-columns inequality gives
$\Omega(L\log d)$ against every code~\cite[Corollary~7.3]{SixteenthPaper2026}. What is not available is the
disjointness.
\item In the abstract model of XOR-frame recursion every level is charged~\cite{FirstPaper2026}, and compilers with
boundedly many stages meet a multiscale obstruction~\cite{ThirdPaper2026}. Both are restricted models.
\end{itemize}
So the logarithm is proved wherever either the coding or the geometry is tied down, and nowhere else.
\Cref{cor:digit,rem:compose-TP} below make one half of this precise in the linear one-hop model: every address digit
must be treated asymmetrically, but one cheap tape operation does that for all digits at once.

\subsection{What the missing statement must be}\label{sec:missing}

The limits above describe the statement that would decide the oblivious case. It must
\begin{itemize}[leftmargin=1.5em]
\item tie what is known to where it can be, for all schedules: the class results of \cref{sec:one-function} do this
for restricted geometries;
\item fail for abstract separated relations, where codes of rate tending to one exist;
\item fail for motion that ignores the data, which padding supplies;
\item constrain all scales at once;
\item and use the shape of tapes, since a bound on coding gains for bounded-degree graphs in general would be a circuit
lower bound (\cref{sec:first}).
\end{itemize}
Since a linear-time multiplier gives a linear-time transposer~\cite{HvdH2025}, a proof of such a statement for all
machines would also be the first superlinear lower bound for off-line multiplication on multitape machines.

\subsection{What a construction cannot avoid}\label{sec:needs}

The same results say what a linear-time machine would have to use.
\begin{itemize}[leftmargin=1.5em]
\item \emph{No bounded complex coefficients.} Linear recovery with bounded complex coefficients and bounded width
stays at the pairing rate on cone relations (\cref{sec:proved}). A linear code would need positive characteristic or
coefficients that cannot be rescaled to be bounded; nonlinear codes are not covered.
\item \emph{Curved motion.} In the limit, head motion must fail to be locally affine on a positive fraction of the
time (\cref{sec:motion}), and a recursive schedule must have a stage whose limiting motion is not regular.
\item \emph{Expanding relations.} The trimmed incidence graphs of the cone relations must not be uniformly hyperfinite
(\cref{sec:proved}). This is necessary and cheap: a fixed linear-time machine has such cones and still obeys the
all-code half bound~\cite[Corollary~7.9]{DirectedRecovery2026}.
\item \emph{No gadgets.} After trimming, almost no output is within bounded distance of its own input bit, so no gain
above the pairing rate can come from gadgets of bounded size (\cref{sec:proved}). The pairing itself is not excluded;
it stops at~$\tfrac12$.
\end{itemize}
Above the pairing rate, at a degree cap that is fixed before the length and the paired distances grow, the only codes
known come from one source: sections of line bundles on the supersingular loci of Hilbert modular
surfaces~\cite{StorageCodes2026}, by a method that originates in~\cite{Luck2026}, together with their
amplifications~\cite{HighRate2026}. At degrees that grow with the length other examples exist.

\subsection{The rule and the relation}\label{sec:rule}

One feature of machines is absent from the coding statements above. A storage code on a relation lets each receiver
choose its own decoding function. A Turing machine has one finite transition table, and on a stretch of time during
which its heads move in step, the function that produces the output at one position is the function that produces the
output at the next position, translated. The cone relation records which inputs each output can depend on. It does not
record that neighbouring outputs depend on them \emph{in the same way}.

Sharing decoders is not by itself an obstruction to rate above one half: by their construction, the linear codes of
the thirteenth paper are equivariant under large groups with boundedly many orbits of coordinates. But transposition has a feature that an
abstract relation lacks: its demand is not invariant under translating one index. Shifting the columns of a matrix
shifts the rows of its transpose. The two theorems of this paper are the two cleanest ways in which structure of the
decoder on one index side is incompatible with the swap: \emph{order} (\cref{sec:order}) and \emph{symmetry}
(\cref{sec:symmetry}). They are proved in the linear one-hop model of \cref{sec:onehop}. They are not theorems about
machines, and \cref{sec:construction} says what would be needed to make them so.

\section{Order: the flag bound}\label{sec:order}

Notation is that of \cref{sec:onehop}: $F$ is a field, $V=F^d$, $\sw$ is the swap, $K=\sum_aA_a\otimes B_a$ with finitely
many terms, and the swap code is $\ker(\sw-K)=\{X:X^{\tr}=\sum_aA_aXB_a^{\tr}\}$. A \emph{complete flag} is a chain of
subspaces $W_1\subset W_2\subset\dots\subset W_d=V$ with $\dim W_i=i$.

\subsection{The theorem}

\begin{theorem}[Theorem~FL, the flag bound]\label{thm:FL}
\partformal{Flag.lean}{finrank_swapCode_le_of_flag,finrank_swapCode_le_of_staircase,finrank_swapCode_le_of_flag_left,finrank_swapCode_swap}{the
choice of bases adapted to the two flags (step~1 of the proof). The Lean statement is the resulting coordinate form:
if $P,Q$ are invertible and every $P^{-1}B_aQ$ has its column $b'$ supported on the first $g(b'+1)$ rows, the bound
holds.}
Let $W_1\subset\dots\subset W_d=V$ be a complete flag and put $g(i)=\dim\sum_aB_a(W_i)$. Then, whatever the $A_a$ are,
\[
  \dim\ker(\sw-K)\ \le\ \sum_{i=1}^{d}g(i).
\]
The same holds with the roles of the $A_a$ and the $B_a$ exchanged.
\end{theorem}

\begin{proof}
\emph{Step 1: bases.} The subspaces $U_i=\sum_aB_a(W_i)$ are nested and $\dim U_i=g(i)$. Choose a basis
$f_0,\dots,f_{d-1}$ of $V$ with $W_i=\spn(f_0,\dots,f_{i-1})$, and a basis $f'_0,\dots,f'_{d-1}$ with
$U_i\subseteq\spn(f'_0,\dots,f'_{g(i)-1})$ for every~$i$; such a basis exists because the $U_i$ are nested: extend a
basis of $U_1$ through $U_2,\dots,U_d$ and then to~$V$. For a term $A\otimes B$ of $K$ write $[B]_{b,b'}$ for the
coefficient of $f'_b$ in $Bf_{b'}$ and $[A]_{a,a'}$ for the coefficient of $f_a$ in $Af'_{a'}$. Since
$Bf_{b'}\in U_{b'+1}$, the coefficient $[B]_{b,b'}$ vanishes unless $b<g(b'+1)$.

\emph{Step 2: the equations.} Write $X=\sum_{a,b}p_{ab}\,f'_a\otimes f_b$. Then $\sw X=\sum p_{ab}\,f_b\otimes f'_a$, whose
coefficient at $f_a\otimes f'_b$ is $p_{ba}$, and the coefficient of $KX$ at $f_a\otimes f'_b$ is
$\sum_{\text{terms}}\sum_{a',b'}[A]_{a,a'}[B]_{b,b'}\,p_{a'b'}$. So $\sw X=KX$ says, for all $a,b$:
\begin{equation}\label{eq:FL}
  p_{ba}\ =\ \text{a linear combination of the }p_{a'b'}\text{ with }g(b'+1)>b .
\end{equation}

\emph{Step 3: elimination.} Let $I=\{(a',b'):a'<g(b'+1)\}$, so $\card I=\sum_ig(i)$. We claim that a solution $p$ that
vanishes on $I$ is zero; then restriction to $I$ is injective on the code and the bound follows. Suppose $p$ vanishes on
$I$, and prove that row $b$ of $p$ vanishes, by downward induction on $b=d-1,\dots,0$. By \eqref{eq:FL}, $p_{ba}$ is a
combination of entries $p_{a'b'}$ with $g(b'+1)>b$. If $a'>b$, row $a'$ vanishes by induction. If $a'\le b$, then
$a'\le b<g(b'+1)$, so $(a',b')\in I$ and $p_{a'b'}=0$. Hence $p_{ba}=0$.

\emph{Step 4: exchange.} Since $\sw(\sw-\sum_aA_a\otimes B_a)\sw=\sw-\sum_aB_a\otimes A_a$, the two kernels have the same
dimension.
\end{proof}

\begin{outline}[Outline of the formal proof of the coordinate form]
In Lean the columns are numbered from~$0$ and the function $g$ is shifted accordingly: Lean's $g(b')$ is the $g(b'+1)$
of the statement above. Let $Q$ be the matrix whose columns are the $f_j$ and $P$ the matrix whose columns are the
$f'_j$; then $(P^{-1}B_aQ)_{b,b'}$ is the coefficient $[B_a]_{b,b'}$. The map $X\mapsto P^{-1}X(Q^{-1})^{\tr}$ is an
isomorphism from the code of $(A_a,B_a)$ onto the code of $(Q^{-1}A_aP,\,P^{-1}B_aQ)$ (\lean{swapCodeEquiv},
\lean{finrank_swapCode_conj}), which reduces the statement to matrices $B_a$ whose column $b'$ vanishes from row
$g(b')$ on. For these, a solution that vanishes on $\{(a',b'):a'<g(b')\}$ is zero, by strong induction on $d-1-b$
(\lean{eq_zero_of_mem_swapCode_of_staircase}), so restriction to that set is injective
(\lean{staircaseRestrict_injective}), and the set has $\sum_{b'}\min(g(b'),d)$ elements (\lean{card_staircase}). The
exchange of sides is \lean{finrank_swapCode_swap}. No monotonicity of $g$ is used.
\end{outline}

The bound depends on one tensor side only, and on that side only through how fast the images of a flag grow: we call
$g$ the \emph{expansion profile} of the $B_a$ along the flag. The transposed profiles are also available, because
$\rank(\sw-K)=\rank(\sw-K^{\tr})$ and $K^{\tr}=\sum_aA_a^{\tr}\otimes B_a^{\tr}$.

\begin{proposition}[The bound is attained for every profile]\label{prop:attained}
\formal{Flag.lean}{finrank_swapCode_staircase,exists_staircase_finrank_eq,mem_swapCode_staircase_iff}
Let $g\colon\{1,\dots,d\}\to\{0,\dots,d\}$ be any function. Take one term for every pair $(i,k)$ with $k<g(i+1)$,
namely $A=E_{ik}$ and $B=E_{ki}$, where $E_{ik}$ is the matrix unit. Then column $i$ of every $B_a$ is supported on
the first $g(i+1)$ rows, and the swap code is $\{X:X_{ki}=0\text{ unless }k<g(i+1)\}$, of dimension exactly
$\sum_ig(i)$.
\end{proposition}

\begin{outline}
Here $KX=\sum X_{ki}E_{ik}$ over the chosen pairs, so $X^{\tr}=KX$ holds at an entry $(i,k)$ with $k<g(i+1)$ and forces
$X_{ki}=0$ at the others (\lean{mem_swapCode_staircase_iff}).
\end{outline}

\begin{remark}[Numerical: the bound is not always tight]\label{rem:not-tight}
For a fixed span of second factors the best flag need not give the true maximum over the first factors. In an
exhaustive computation over $\F_2$ at $d=3$, for $4$ of the $20$ orbits of spans of at most two second factors the
maximum is $4$ while the best flag, over $\F_2$, $\F_4$ or $\F_8$, gives $5$ or $6$.
\end{remark}

\subsection{Corollaries}

Write $h=d(d-1)/2$, so that $\dim\ker(\sw-K)\le d(d+1)/2$ is the same as $\rank(\sw-K)\ge h$.

\begin{corollary}[A pair of flags]\label{cor:triangular}
\partformal{Flag.lean}{two_mul_finrank_swapCode_le_of_triangular,two_mul_finrank_swapCode_le_of_triangular_left}{the
passage from a pair of flags to bases in which the $B_a$ are upper triangular, as in step~1 of the proof of
\cref{thm:FL}. The Lean statement is the matrix form.}
If $P^{-1}B_aQ$ is upper triangular for all $a$, for some invertible $P$ and~$Q$, then $\dim\ker(\sw-K)\le d(d+1)/2$,
for arbitrary $A_a$. In terms of flags: this holds if there are complete flags $W$ and $W'$ with
$B_a(W_i)\subseteq W'_i$ for all $a$ and~$i$.
\end{corollary}

\begin{outline}
This is \cref{thm:FL} with $g(i)\le i$, and $\sum_{i\le d}i=d(d+1)/2$ (\lean{two_mul_sum_fin_succ}). Bases adapted to
$W$ and $W'$ make the $B_a$ upper triangular, and conversely.
\end{outline}

So the pairing value is the value of a hidden order: a pair of flags is an order in which the $t$-th generalized column
is computed from the first $t$ generalized rows.

\begin{corollary}[At most two terms, every field]\label{cor:two-terms}
\notformalized
For all $A_1,A_2,B_1,B_2\in\Mat_d(F)$,
\[
  \dim\{X:X^{\tr}=A_1XB_1^{\tr}+A_2XB_2^{\tr}\}\ \le\ \frac{d(d+1)}2 .
\]
\end{corollary}

\begin{informalproof}
The dimension of a kernel does not change under extension of the field, so let $F$ be algebraically closed. We show
that any two linear maps $B_1,B_2\colon V\to V'$ between spaces of the same dimension $d$ admit complete flags with
$B_j(W_i)\subseteq W'_i$; \cref{cor:triangular} then applies. The determinant $\det(\alpha B_1+\beta B_2)$, in any bases,
is a binary form of degree~$d$, so it vanishes at some $(\alpha,\beta)\ne0$ (at every $(\alpha,\beta)$ if it is
identically zero). Take $v\ne0$ with $(\alpha B_1+\beta B_2)v=0$. Then $B_1v$ and $B_2v$ are proportional; let
$W_1=\langle v\rangle$ and let $W'_1$ be a line containing both. The maps induced on $V/W_1\to V'/W'_1$ are again two
maps between spaces of equal dimension, and induction on~$d$ gives the flags.
\end{informalproof}

\begin{remark}[Two is sharp]\label{rem:three}
Three terms can exceed the pairing value. For every $d$, the operator $\Pi=\sum_{i,j}E_{ij}\otimes E_{ij}$ is
$X\mapsto\operatorname{tr}(X)\cdot1$, of rank one, and
\[
  \sw-\Pi\ =\ \sum_{i<j}(E_{ij}-E_{ji})\otimes(E_{ji}-E_{ij})
\]
has $h$ terms. So $K=\sw-\Pi$ has $h$ terms and its code $\ker\Pi$, the matrices of trace zero, has dimension $d^2-1$,
in every characteristic. At $d=3$ that is three terms and dimension $8>6$. At $d=4$, adding $E_{33}$ to both factors
of the third term of the $d=3$ example, with indices $0,\dots,3$, gives three terms and a code of dimension $11>10$
(numerical: checked over $\F_2,\F_3,\F_5,\F_7$ by \leanfile{tools/three_term_example.py}). We have no three-term
example for $d\ge5$.
\end{remark}

\begin{corollary}[One triangularizable side]\label{cor:commutative}
\notformalized
Suppose that, over the algebraic closure of $F$, the $B_a$ have a common invariant complete flag. Then
$\dim\ker(\sw-K)\le d(d+1)/2$ for any number of terms and arbitrary $A_a$. This holds in particular when
\begin{enumerate}[label=(\alph*),leftmargin=2em]
\item the $B_a$ commute pairwise, for example when all are circulant;
\item $\chr F=p$ and all $B_a$ lie in the linear span of one finite $p$-group of matrices.
\end{enumerate}
\end{corollary}

\begin{informalproof}
An invariant flag gives \cref{cor:triangular} with $W'=W$, over the closure, and the dimension does not depend on the
field. (a) Commuting matrices over an algebraically closed field have a common eigenvector: an eigenspace of one of them
is invariant under the others, and one inducts on the dimension. Passing to the quotient by that eigenvector and
inducting again gives a common invariant flag. (b) Let $G$ be a finite $p$-group of invertible matrices in
characteristic~$p$ and $z\ne1$ a central element of order~$p$. Then $(z-1)^p=z^p-1=0$, so $\ker(z-1)$ is a nonzero
$G$-invariant subspace on which $G/\langle z\rangle$ acts; by induction on the order of the group, $G$ fixes a nonzero
vector. Passing to the quotient by that vector and inducting on the dimension gives a complete flag invariant under
$G$, hence under its span.
\end{informalproof}

For circulant $B_a$ the statement reads: \emph{if every column of $X$ is a linear function of the rows of $X$ that
commutes with cyclic rotation of the rows' coordinates (any number of rows, each rotated and convolved with kernels that
may depend on the column), then $X$ ranges over a space of dimension at most $d(d+1)/2$.} Reading any number of source
rows in step with the output cannot beat the symmetric matrices.

In (b) the group must be the same for all $B_a$. For $p=2$ and $d=2^m$ two natural examples are the automorphisms of
the hierarchy of dyadic intervals (block swaps, block reversals, XOR translations) and the affine group
$\{i\mapsto ui+t:u\text{ odd}\}$ of $\Z/2^m$. They are different subgroups of the symmetric group, and mixing them
is not covered: at $d=4$ a rotation and the transposition $(0\;1)$, which swaps the halves of one dyadic block of length
two, generate the whole symmetric group. The next corollary
covers more of the second example.

\begin{corollary}[All affine index maps in characteristic $p$]\label{cor:affine}
\notformalized
Let $\chr F=p$ and $d=p^m$. For $u,t\in\Z/d$ let $\varphi_{u,t}$ be the linear map on $F^{\Z/d}$ with
$(\varphi_{u,t}x)_i=x_{ui+t}$; here $u$ need not be a unit. If every $B_a$ lies in the span of the $\varphi_{u,t}$, then
$\dim\ker(\sw-K)\le d(d+1)/2$.
\end{corollary}

\begin{informalproof}
Let $R=\varphi_{1,1}$ be the rotation. Then $R\varphi_{u,t}=\varphi_{u,t}R^u$, since both send $x$ to
$(x_{ui+t+u})_i$. Hence $(R-1)^i\varphi_{u,t}=\varphi_{u,t}(R^u-1)^i$, and $R^u-1$ is a multiple of $R-1$ in the
commutative algebra $F[R]$, so $\varphi_{u,t}$ maps $\ker(R-1)^i$ into itself. As $R$ is a single cycle of length
$p^m$ in characteristic~$p$, $R-1$ is nilpotent with a one-dimensional kernel, and the subspaces $\ker(R-1)^i$,
$1\le i\le d$, form a complete flag. Apply \cref{cor:triangular} with $W'=W$.
\end{informalproof}

So over $\F_2$, on rows of length $2^m$, rotations, reflections, dilations, wrap-around decimations and convolutions
all preserve one flag, in any combination and at any depth (\cref{rem:composition}). What leaves that world is a
projection: a mask that selects part of a row.

\begin{corollary}[Blocks and expansion]\label{cor:expansion}
\notformalized
\begin{enumerate}[label=(\alph*),leftmargin=2em]
\item If, for suitable invertible $P,Q$, all $P^{-1}B_aQ$ are block upper triangular with diagonal blocks of sizes
$d_1,\dots,d_m$, then $\dim\ker(\sw-K)\le(d^2+\sum_sd_s^2)/2$.
\item If some flag has $g(i)\le ci+c_0$ with $c\ge1$ and $c_0\ge0$, then
$\dim\ker(\sw-K)\le\bigl(1-\frac1{2c}\bigr)d^2+\bigl(\frac{c_0}c+\frac12\bigr)d$.
\item A sum of $k$ terms has $\dim\ker(\sw-K)\le\sum_{i\le d}\min(d,ki)\le\bigl(1-\frac1{2k}\bigr)d^2+\frac d2$.
\item Let $d$ be prime to $\chr F$ and $s\ge1$. If every $B_a$ lies in the span of the maps $\varphi_{u,t}$ of
\cref{cor:affine} with $u\in\{-s,\dots,s\}$, then $\dim\ker(\sw-K)\le\bigl(1-\frac1{2s}\bigr)d^2+\bigl(\frac1s+\frac12\bigr)d$.
\end{enumerate}
\end{corollary}

\begin{informalproof}
(a) Take $g(i)=d_1+\dots+d_s$ for $i$ in block~$s$; then $\sum_ig(i)=\sum_sd_s(d_1+\dots+d_s)=(d^2+\sum_sd_s^2)/2$.
(b) If $d<c_0$ the right side exceeds $d^2$. Otherwise put $u=(d-c_0)/c$. Then
$\sum_i\min(d,ci+c_0)=d^2-\sum_{i\le u}c(u-i)\le d^2-\frac c2(u^2-u)$, because
$\sum_{i\le\lfloor u\rfloor}(u-i)\ge(u^2-u)/2$, and $d^2-\frac{(d-c_0)^2}{2c}+\frac{d-c_0}2$ is at most the claimed
bound. (c) Any flag has $g(i)\le ki$; apply~(b). (d) Extend the field so that it contains a primitive $d$-th root of
unity, and let $\chi_k$ be the characters of $\Z/d$. Then $\varphi_{u,t}\chi_k$ is a multiple of $\chi_{uk}$. Order the
frequencies $0,1,-1,2,-2,\dots$ and let $W_i$ be spanned by the first $i$ characters. The frequencies $uk$ with
$\abs u\le s$ and $\chi_k\in W_i$ lie in an interval of at most $si+1$ residues, so $g(i)\le si+1$; apply~(b).
\end{informalproof}

The hypothesis $c\ge1$ in (b) is needed: for $B=\mathrm{diag}(1_{d/2},0)$, $A=1$ and a flag that alternates kernel
and non-kernel vectors, $g(i)=\lfloor i/2\rfloor$ and the code has dimension $d(d+2)/8$.

\begin{remark}[Entries of $b$ bits]\label{rem:entries}
For $K=\sum_a\tilde A_a\otimes B_a$ on $V\otimes V\otimes F^b$ with $\tilde A_a\in\operatorname{End}(V\otimes F^b)$
arbitrary and $B_a\in\operatorname{End}(V)$, the proof of \cref{thm:FL} gives
$\dim\ker(\sw\otimes1-K)\le b\sum_ig(i)$, with coefficients $p_{ab}\in F^b$. Operators that mix the column index
with the position inside an entry are covered as well: every operator on $V\otimes F^b$ is a sum $\sum_jB_j\otimes C_j$,
the $C_j$ can be absorbed into the first factors, and the profile is that of the $B_j$. For a cyclic shift of the bit
string of a row by a number of bits that is not a multiple of~$b$, the $B_j$ are two rotations.
\end{remark}

\begin{remark}[Composition]\label{rem:composition}
If the second factors of $K_1,\dots,K_t$ all lie in one algebra $\mathfrak B\subseteq\Mat_d(F)$, so do those of the
product $K_t\cdots K_1$, because $(A\otimes B)(A'\otimes B')=AA'\otimes BB'$. So \cref{cor:commutative,cor:affine} are
statements about linear data paths of any depth: a path that acts on one index through a commutative algebra, or through
the affine maps of \cref{cor:affine}, and on the other index in any way at all, is correct on at most the pairing value.
\end{remark}

\subsection{What the theorem says}\label{sec:order-says}

\begin{itemize}[leftmargin=1.5em]
\item \emph{Only the weaker side counts.} The bound uses the expansion of the $A_a$ or of the $B_a$, whichever is
smaller. Arbitrary mixing on one index buys nothing while the other index has a flag.
\item \emph{A code above the pairing value needs expansion on both sides}: for every flag on either side, the images of
the flag must grow faster than the flag. By \cref{cor:expansion}(b), a profile $g(i)\le ci$ with $c\ge1$ gives rate at
most $1-1/(2c)+O(1/d)$, so rate $1-\eps$ needs expansion factor about $1/(2\eps)$ along every flag, on both sides.
\item \emph{Expansion is cheap on a tape.} Over the algebraic closure, rotations of a row have an invariant flag, and
so do masks by intervals, but the two flags are incompatible: rotations together with the mask of one proper interval
generate the full matrix algebra. For rotations with the masks of $M$ intervals $J_l$, a rotation-invariant flag has
$g(i)=\sum_l\min(i,\card{J_l})\le Mi$, and no flag does better, because the rotations of $W_i$ span a
rotation-invariant space of dimension at least~$i$, and such a space of dimension $h$ restricts to an interval of
length $m$ in dimension $\min(h,m)$. For $M$ equal intervals this gives rate $1-1/(2M)$ and no more. So the flag
bound cannot hold that family to the pairing rate. The next section does.
\end{itemize}

\section{Symmetry: the translate pigeonhole}\label{sec:symmetry}

The flag bound cannot hold rotations with interval masks to the pairing rate (\cref{sec:order-says}). This section
does, by a different argument: a code larger than half the space meets its own translate, and on the intersection a
shared decoder forces the data to be translation-covariant.

\subsection{The lemma}

For operators $S,T$ write $[S,T]=ST-TS$.

\begin{lemma}[Pigeonhole against a translate]\label{lem:pigeonhole}
\formal{Pigeonhole.lean}{two_mul_finrank_ker_sub_le_ker,finrank_ker_commutator_sub_le,two_mul_finrank_ker_sub_le,two_mul_finrank_ker_sub_le_of_commute,finrank_range_commutator_le}
Let $\sw,K,T$ be linear operators on a space of finite dimension $N$ over a field. Then
\[
  \dim\ker(\sw-K)\ \le\ \tfrac12\bigl(N+\dim\ker[\sw-K,T]\bigr)\ \le\ \tfrac12\bigl(N+\dim\ker[\sw,T]+\rank[K,T]\bigr).
\]
\end{lemma}

\begin{proof}
Let $\X=\ker(\sw-K)$, $D=\dim\X$ and $\X_T=\{X\in\X:TX\in\X\}$. The map $\X\to F^N/\X$, $X\mapsto TX+\X$, has kernel
$\X_T$ and a target of dimension $N-D$, so $\dim\X_T\ge2D-N$; $T$ need not be invertible. For $X\in\X_T$,
$(\sw-K)TX=0=T(\sw-K)X$, so $\X_T\subseteq\ker[\sw-K,T]$. That is the first inequality. For the second, $[\sw-K,T]X=0$
means $[\sw,T]X=[K,T]X\in\im[K,T]$, and the preimage of a subspace $W$ under a map has dimension at most that of its
kernel plus $\dim W$.
\end{proof}

With $S=\sw-K$ the first inequality says $\rank[S,T]\le2\rank S$: a commutator has rank at most twice the rank of
either factor. The Lean proof takes this route, from $\rank[S,T]\le\rank(ST)+\rank(TS)$ and rank--nullity
(\lean{finrank_range_commutator_le}). The lemma is elementary, and we claim no novelty for it. What is used below is the form with $\sw$ and $K$ separated: $K$
may nearly commute with a $T$ with which $\sw$ does not.

\begin{proposition}[For the swap, the best case is the pairing value]\label{prop:swap-floor}
\partformal{Swap.lean}{card_le_finrank_ker_swap_commutator,two_mul_finrank_ker_swap_sub_one,two_mul_finrank_range_swap_sub_one}{the
equality case for $T=1\otimes G$.}
On $V\otimes V$ with $\dim V=d$, every operator $T$ has $\dim\ker[\sw,T]\ge d$, in every characteristic. Equality holds
for $T=1\otimes G$, that is $X\mapsto XG^{\tr}$, when $G$ is non-derogatory.
\end{proposition}

\begin{proof}
The kernel of $\sw-1$ is the space of symmetric matrices, of dimension $d(d+1)/2$ in every characteristic, so
$\rank(\sw-1)=d(d-1)/2$, and $[\sw,T]=[\sw-1,T]$ has rank at most $d(d-1)$ by \cref{lem:pigeonhole}. For $T=1\otimes G$,
$[\sw,T]X=GX^{\tr}-X^{\tr}G^{\tr}$, so the kernel is the space of $Y=X^{\tr}$ with $GY=YG^{\tr}$. As $G^{\tr}$ is
similar to $G$, this space is isomorphic to the centralizer of $G$, whose dimension is $d$ exactly when $G$ is
non-derogatory.
\end{proof}

So \cref{lem:pigeonhole} never gives less than $(d^2+d)/2$ for the swap, and it gives exactly the pairing value when
$K$ commutes with a suitable~$T$. The rest of the section is about $K$ that commute with a translation up to a
boundary.

\subsection{Class-uniform rules}

Let $\sigma$ be the cyclic rotation of $F^{\Z/d}$, $(\sigma x)_j=x_{j+1}$. For a matrix $X$ with rows
$x_c\in F^{\Z/d}$, $c\in\Z/d$, write $\sigma^\delta X$ for the matrix with rows $\sigma^\delta x_c$.

\begin{definition}[Class-uniform rules]\label{def:class-uniform}
\formal{ClassUniform.lean}{rot,IsClassUniformCode}
Fix a partition of the row index $\Z/d$ into non-empty classes. A \emph{class-uniform rule} is a family of linear maps
$\cA_J\colon\Mat_d(F)\to F^{\Z/d}$, one for each class~$J$. A subspace $\X\subseteq\Mat_d(F)$ is a \emph{code} of the
rule if
\[
  x_c=\cA_J(\sigma^cX)\qquad\text{for every }X\in\X,\text{ every class }J\text{ and every }c\in J .
\]
\end{definition}

In words: row $c$ of $X$ is a fixed linear function of the whole matrix with its columns rotated by~$c$, and the function
depends only on the class of~$c$. The receiver of row $c$ looks at the stored word from its own position, through a
window of any size, and receivers in one class use the same decoder.

In the notation of \cref{sec:onehop} (this translation is not formalized): let $K=\sum_aA_a\otimes B_a$ have second
factors in the span of the matrices $P_J\sigma^e$, a rotation followed by the mask $P_J$ of a class, and arbitrary first
factors, say $K=\sum_{J,e}A_{J,e}\otimes P_J\sigma^e$. Then $(KX)_{r,c}=\sum_{i,e}(A_{J,e})_{r,i}X_{i,c+e}$ for $c\in J$,
so the swap code of $K$ is the largest code of the rule $\cA_J(Y)_r=\sum_{i,e}(A_{J,e})_{r,i}Y_{i,e}$, and every rule
arises in this way. This is the family of \cref{sec:order-says}. For a single class the $B_a$ are circulant and
\cref{cor:commutative} applies. If the mask is applied before the rotation, $K^{\tr}$ has second factors
$P_J\sigma^{-e}$, and since $\rank(\sw-K)=\rank(\sw-K^{\tr})$ the bounds below hold for that family as well.

\begin{theorem}[Theorem~TP, the translate pigeonhole]\label{thm:TP}
\formal{ClassUniform.lean}{IsClassUniformCode.finrank_le,Link,IsClassUniformCode.two_mul_finrank_le,IsClassUniformCode.two_mul_finrank_le_of_const}
Let $\X$ be a code of a class-uniform rule. Let $\Delta_0\subseteq\Z/d$ be a set of $s+1\ge2$ shifts, and let $\chi$
be the number of connected components of the graph on $\Z/d$ in which $c\sim c'$ when $c$ and $c'$ lie in one class and
$c'-c\in\Delta_0-\Delta_0$. Then
\[
  \dim\X\ \le\ \frac{s\,d^2+d\,\chi}{s+1}.
\]
In particular, if the classes are $M$ cyclic intervals, then $\dim\X\le(d^2+Md)/2$.
\end{theorem}

\begin{outline}
Let $\X'=\{X:\sigma^\delta X\in\X\text{ for all }\delta\in\Delta_0\}$, an intersection of $s+1$ preimages of~$\X$;
each has codimension at most $d^2-\dim\X$, so $\dim\X'\ge(s+1)\dim\X-s\,d^2$
(\lean{card_mul_finrank_add_le_iInf_comap}, \lean{mul_finrank_le_add_finrank_iInf_comap}, in
\leanfile{Pigeonhole.lean}). Let $X\in\X'$, $\delta,\delta'\in\Delta_0$ and $c,c'$ in one class $J$
with $c+\delta=c'+\delta'$. The rule applied to $\sigma^\delta X$ at row $c$ gives
$\sigma^\delta x_c=\cA_J(\sigma^{c+\delta}X)$, and applied to $\sigma^{\delta'}X$ at row $c'$ it gives
$\sigma^{\delta'}x_{c'}=\cA_J(\sigma^{c'+\delta'}X)$. The right sides agree, so $x_{c'}=\sigma^{c'-c}x_c$
(\lean{IsClassUniformCode.row_eq_of_rot_mem}). Hence each row of $X$ determines the rows in its component, restriction
to one row per component is injective on $\X'$, and $\dim\X'\le d\chi$. The Lean statement takes any set $S$ of rows
that meets every component and gives $(s+1)\dim\X\le s\,d^2+d\card S$. For $M$ cyclic intervals and $\Delta_0=\{0,1\}$ the components are the intervals. In Lean this case
(\lean{IsClassUniformCode.two_mul_finrank_le}) takes a non-empty set $S$ containing every row whose class differs from
that of the row before it, so that $\card S=M$ for $M\ge2$ intervals and $S$ is any single row for one class
(\lean{IsClassUniformCode.two_mul_finrank_le_of_const}).
\end{outline}

The space on which the proof lands is worth naming. For one unit shift it is
$\Hk=\{X:x_{c+1}=\sigma x_c\text{ whenever }c,c+1\text{ lie in one class}\}$: inside each interval the matrix is constant
along anti-diagonals. It has $d$ free parameters per interval, whatever the length of the interval.

We say that a partition into $M$ classes has an \emph{almost-period} if some single shift $\delta$ leaves only $M$
components in the graph of \cref{thm:TP} with $\Delta_0=\{0,\delta\}$, so that the bound is $(d^2+Md)/2$. Intervals have
the almost-period~$1$, and residue classes modulo $M$ have the almost-period~$M$.

\begin{lemma}[Extremal codes contain $\Hk$]\label{lem:extremal}
\notformalized
Let the classes be $M$ cyclic intervals and let $\X$ be a code of a class-uniform rule with $2\dim\X=d^2+Md$. Then
$\X\cap\sigma^{-1}\X=\Hk$; in particular $\Hk\subseteq\X$.
\end{lemma}

\begin{informalproof}
By the proof of \cref{thm:TP}, $\X'=\X\cap\sigma^{-1}\X$ lies in $\Hk$ and has dimension at least
$2\dim\X-d^2=Md=\dim\Hk$.
\end{informalproof}

\begin{proposition}[Equal pieces: the bound is attained]\label{prop:block-symmetric}
\formal{BlockSymmetric.lean}{blockSymmetric,isClassUniformCode_blockSymmetric,two_mul_finrank_blockSymmetric,blockSymmetric_attains_bound}
Let $d=Mm$ with $M,m\ge1$ and let the classes be the $M$ intervals of length~$m$. The \emph{block-symmetric} matrices,
those whose $m\times m$ blocks in the $M\times M$ grid are all symmetric, are a code of a class-uniform rule, of
dimension $(d^2+Md)/2$.
\end{proposition}

\begin{outline}
Write a row index as $c=sm+k$ and a column index as $r=tm+l$ with $0\le k,l<m$. Block symmetry says
$X_{sm+k,\,tm+l}=X_{sm+l,\,tm+k}$. The right side is the entry of $\sigma^cX$ in row $sm+l$ and column
$tm+k-c=(t-s)m$, which depends on $c$ only through its block~$s$. So $x_c=\cA_s(\sigma^cX)$ with
$\cA_s(Y)_{tm+l}=Y_{sm+l,\,(t-s)m}$ (\lean{isClassUniformCode_blockSymmetric}). There are $M^2$ blocks, each with
$m(m+1)/2$ free entries, and $M^2m(m+1)/2=(d^2+Md)/2$ (\lean{two_mul_finrank_blockSymmetric}); the set of the $M$
first rows of the blocks satisfies the hypothesis of \cref{thm:TP} for the unit shift
(\lean{blockSymmetric_attains_bound}).
\end{outline}

The operator $K$ of this rule is the coarse transposition of the grid: $X\mapsto(\cA_{s(c)}(\sigma^cX))_c$, read as a
matrix with these columns, has the entry $X_{sm+l,\,tm+k}$ in position $(tm+l,\,sm+k)$, so block $(s,t)$ is moved to
position $(t,s)$ with its inner coordinates kept.

So the largest code of a class-uniform rule with $M$ equal interval pieces has dimension exactly $(d^2+Md)/2$, the
value of the $M\times M$ coarse transposition, against the bound $1-1/(2M)$ that flags give for the same family. For
$M=d$ the pieces are single rows, the rule is unrestricted and the code is everything. In between, each additional
piece raises the bound by $d/2$.

\begin{remark}[Numerical: what the exact value does not say]\label{rem:extremal}
Exhaustive computations over $\F_2$ show the following.
(a) Extremal codes are far from unique: for $d=4$ and classes $\{0,1\},\{2,3\}$ there are $59\,392$ codes of
dimension~$12$, in $327$ orbits under the evident symmetries, and the block-symmetric code is one orbit.
(b) Coding beats copying inside the bound: for $d=4$ and pieces $\{0\},\{1,2,3\}$ the bound $12$ is attained, but no
rule that copies one stored entry to each output entry exceeds~$11$.
(c) The bound is not attained for every partition into unequal intervals: by \cref{lem:extremal} and a solver, the
pieces $\{0\},\{1\},\{2,3,4\}$ at $d=5$ have largest code $19<20$ (\cref{rem:small}).
\end{remark}

\begin{proposition}[Copy rules, every partition]\label{prop:copy}
\notformalized
Call a class-uniform rule a \emph{copy rule} if every output entry is a multiple of one stored entry:
$X_{c,r}=\lambda_{J,r}\,X_{\pi_J(r),\,\rho_J(r)+c}$ for $c\in J$. A code of a copy rule with $M$ classes, of any shape,
has dimension at most $(d^2+Md)/2$.
\end{proposition}

\begin{informalproof}
Let $f(c,r)=(\pi_J(r),\rho_J(r)+c)$ for $c\in J$, a map on the $d^2$ positions. A matrix in the code is determined, on
each component of the functional graph of $f$, by its value at one position on the cycle of that component, so the
dimension is at most the number of cycles. A fixed point $(c,r)$ has $\pi_J(r)=c$ with $J$ the class of $c$, so for each
$r$ there is at most one fixed point per class: at most $Md$ in all. Every other cycle has at least two positions.
\end{informalproof}

So for copy rules the shape of the classes does not matter, while for coded rules \cref{thm:TP} needs the classes to
have an almost-period. The difference is real: \cref{prop:aperiodic} exhibits two classes with no almost-period and a
coded rule whose code has dimension greater than $(d^2+Md)/2$. How far coded rules on such classes can go is
\cref{prob:aperiodic}.

\subsection{On a line}

A tape is a line, not a cycle. Let rows and columns be indexed by $\{0,\dots,d-1\}$ and let the classes be intervals. A
\emph{two-sided rule on a line} is given by coefficients $\cA_{J,i}[r][e]$, $-d<e<d$, with
\[
  X_{c,r}=\sum_{i}\sum_{j=0}^{d-1}\cA_{J,i}[r][j-c]\,X_{i,j}\qquad(c\in J).
\]
The decoder of row $c$ again sees the stored word from position $c$, through a window that is cut off at the ends
instead of wrapping around. Every class-uniform rule is of this form, with the same classes: take
$\cA_{J,i}[r][e]$ to depend on $e$ modulo~$d$ (this remark is not formalized). On a line, an interval is a set of
consecutive rows that does not wrap around.

\begin{theorem}[The line form of Theorem~TP]\label{thm:line}
\formal{Line.lean}{IsLineUniformCode,IsLineUniformCode.two_mul_finrank_le,isLineUniformCode_of_coeff,lineView,truncShift}
A code of a two-sided rule on a line with $M$ interval classes has dimension at most $(d^2+Md)/2$.
\end{theorem}

\begin{outline}
Let $(TX)_{i,j}=X_{i,j+1}$ for $j<d-1$ and $(TX)_{i,d-1}=0$, and $\X'=\{X\in\X:TX\in\X\}$, so that
$\dim\X'\ge2\dim\X-d^2$. Let $X\in\X'$ and let $c,c+1$ lie in one class~$J$. The rule for $TX$ at row $c$ reads
$X_{c,r+1}=\sum_i\sum_{j\ge1}\cA_{J,i}[r][j-(c+1)]X_{i,j}$, with $X_{c,d}$ read as~$0$, and the rule for $X$ at row
$c+1$ has the same sum with $j\ge0$. So
\[
  X_{c+1,r}-X_{c,r+1}=\sum_i\cA_{J,i}[r][-(c+1)]\,X_{i,0}.
\]
Suppose that column $0$ of $X$ vanishes. Then inside each class, with rows $a,\dots,a+m_J-1$, the matrix is constant
along anti-diagonals. An anti-diagonal that meets column~$0$ inside the class is zero, and so is one whose first
position in the class is in the last column below row~$a$, because $X_{c,d}=0$. The remaining anti-diagonals start in
row~$a$ at a column $r\ge m_J$: there are $d-m_J$ of them. So
$\dim\X'\le d+\sum_J(d-m_J)=Md$. In Lean the decoder of a class is an arbitrary linear function of the zero-padded
window $(X_{i,e+c})_{i,e}$ (\lean{lineView}); the coefficient form above is an instance
(\lean{isLineUniformCode_of_coeff}). The classes need not be intervals there: the bound is $(d^2+d\card S)/2$ for any
set $S$ of rows that contains row~$0$ and every row whose class differs from that of the row before it. The key step is
\lean{IsLineUniformCode.row_succ_eq}.
\end{outline}

With one class this gives the pairing value for second factors in the span of the truncated shifts, the Toeplitz
matrices, which do not commute and are not covered by \cref{cor:commutative}; the symmetric matrices attain it
(\lean{isLineUniformCode_symmetric}). For $M$ equal intervals the block-symmetric matrices attain the bound on a line
as well, by the remark above; this is not formalized.

\begin{remark}[Numerical: on a line, other shifts fail]\label{rem:line-fails}
For $d=3$ and classes $\{0,2\},\{1\}$, \cref{thm:TP} with $\Delta_0=\{0,2\}$ gives $7$ on a cycle. On a line the largest
code of a two-sided rule has dimension~$8$ (exhaustive over $\F_2$); one such code is $\{X:X_{0,0}=X_{2,2}\}$. So
\cref{thm:line} does not extend to shift sets other than $\{0,1\}$.
\end{remark}

\subsection{Twists, digits and composition}

A rule need not commute with a translation to be caught; it is enough that it intertwines it with boundedly many others.

\begin{theorem}[Twist classes]\label{thm:twist}
\formal{Twist.lean}{succ_mul_finrank_swapCode_le_of_twist_classes,two_mul_finrank_swapCode_le_of_twist,twistSpace,succ_mul_finrank_swapCode_le,finrank_le_sum_finrank_twistSpace}
Let $G\in\Mat_d(F)$, let the terms of $K=\sum_aA_a\otimes B_a$ fall into classes $1,\dots,s$, and let
$G'_1,\dots,G'_s$ commute pairwise, with $GB_a^{\tr}=B_a^{\tr}G'_\kappa$ for every $a$ in class~$\kappa$. Then
\[
  \dim\ker(\sw-K)\ \le\ \frac{s\,d^2+\sum_\kappa h_\kappa}{s+1},\qquad h_\kappa=\dim\{P:G^{\tr}P=PG'_\kappa\}.
\]
\end{theorem}

\begin{outline}
Let $\X'=\{X:XG^t\in\ker(\sw-K),\ 0\le t\le s\}$, of dimension at least $(s+1)\dim\ker(\sw-K)-s\,d^2$
(\lean{succ_mul_finrank_swapCode_le}). For $X\in\X'$ and $P=X^{\tr}$ the hypothesis gives
$(G^{\tr})^tP=\sum_\kappa T_\kappa G_\kappa'^{\,t}$ for $0\le t\le s$, with $T_\kappa=\sum_{a\in\kappa}A_aXB_a^{\tr}$. Applying
$P\mapsto G^{\tr}P-PG'_\kappa$ removes class~$\kappa$ from such a system and lowers $s$ by one, and its kernel is the
space counted by $h_\kappa$; induction on the number of classes gives $\dim\X'\le\sum_\kappa h_\kappa$
(\lean{finrank_le_sum_finrank_twistSpace}). This is the statement that $\prod_\kappa(L-R_\kappa)P=0$ for the commuting
operators $L$, left multiplication by $G^{\tr}$, and $R_\kappa$, right multiplication by $G'_\kappa$.
\end{outline}

\begin{corollary}[Modulation frequencies]\label{cor:weyl}
\notformalized
Let $F$ contain a primitive $d$-th root of unity $\omega$ and let $Z=\mathrm{diag}(1,\omega,\dots,\omega^{d-1})$. If
every $B_a^{\tr}$ is a product $Z^kC$ with $C$ circulant and $k$ in a fixed set of $s$ frequencies, then
$\dim\ker(\sw-K)\le\frac s{s+1}(d^2+d)$.
\end{corollary}

\begin{informalproof}
$\sigma Z=\omega Z\sigma$, so $\sigma Z^kC=Z^kC\,(\omega^k\sigma)$: take $G=\sigma$ and $G'_k=\omega^k\sigma$, which
commute. The matrices $\sigma^{\tr}=\sigma^{-1}$ and $\omega^k\sigma$ are diagonalizable with the same simple spectrum,
the $d$-th roots of unity, so $h_k=d$.
\end{informalproof}

One frequency is the pairing rate, and all $d$ frequencies span the full matrix algebra. The mask of an interval
contains many frequencies (all of them when its length is prime to~$d$), which is why interval masks needed
\cref{thm:TP}.

\begin{corollary}[Every address digit must be broken]\label{cor:digit}
\partformal{Pigeonhole.lean}{two_mul_finrank_ker_sub_le_of_commute}{everything specific to the statement: that $K$
commutes with~$T$, the dimension count $\dim\ker[\sw,T]=d^2/q$, and the sharpness. Lean has only the abstract
lemma.}
Write the index set as $\Z/m\times\Z/q$, $d=mq$, so that the second coordinate is one digit of the address. Suppose
every $B_a$ has the form $B'_a\otimes C_a$ with all $C_a$ in $F[G]$ for one non-derogatory $q\times q$ matrix $G$, for
example with all $C_a$ circulant. Then
\[
  \dim\ker(\sw-K)\ \le\ \tfrac12\Bigl(d^2+\frac{d^2}q\Bigr),
\]
with equality for the partial transpose that exchanges the first digits of the two indices and leaves the second
digits in place.
\end{corollary}

\begin{informalproof}
$K$ commutes with $T=1\otimes(1_m\otimes G)$, and as in \cref{prop:swap-floor} $\ker[\sw,T]$ is isomorphic to the
centralizer of $1_m\otimes G$, which is $\Mat_m(F)\otimes F[G]$, of dimension $m^2q=d^2/q$. Apply
\cref{lem:pigeonhole}. The partial transpose is $\sw_m\otimes1$ on $(F^m\otimes F^m)\otimes(F^q\otimes F^q)$; its second
factors are $E_{u'u}\otimes1_q$, and its swap code is $\Mat_m(F)\otimes\{\text{symmetric }q\times q\text{ matrices}\}$,
of dimension $m^2q(q+1)/2$.
\end{informalproof}

The transposition is the product of the commuting exchanges of the digits of the two addresses
(\cref{sec:saturation}). \Cref{cor:digit} says that a data path which treats any one digit of one address uniformly
is at rate $\tfrac12+\tfrac1{2q}$, however it treats the others; by \cref{lem:pigeonhole}, rate $1-\eps$ needs
$\rank[K,T]\ge(1-2\eps-\tfrac1q)d^2$ for the rotation $T$ of every digit of either index. This is a requirement at
every scale.

\begin{remark}[Composition, and why this gives no logarithm]\label{rem:compose-TP}
\Cref{thm:TP} is a statement about one hop: the class-uniform rules with $M$ classes are not closed under
composition, and two of them compose to a rule with up to $d$ classes. \Cref{lem:pigeonhole} does compose. Since
$[K_2K_1,T]=K_2[K_1,T]+[K_2,T]K_1$, $\rank[K_2K_1,T]\le\rank[K_1,T]+\rank[K_2,T]$
(\lean{finrank_range_commutator_comp_le}), and by induction
\[
  \rank[K_t\cdots K_1,T]\ \le\ \sum_i\rank[K_i,T].
\]
For $T=1\otimes\sigma$ and a hop whose second factors lie in the span of the $P_J\sigma^e$ with $M$ interval classes,
$[K_i,T]X$ is supported on the $M$ columns at which the class changes, so $\rank[K_i,T]\le Md$ and $t$ such hops give
$\dim\ker(\sw-K_t\cdots K_1)\le(d^2+d+tMd)/2$.

The same subadditivity shows why \cref{cor:digit} does not produce a factor $\log d$. An exact transposer needs
$\sum_i\rank[K_i,T]\ge d^2-d^2/q$ for the rotation $T$ of every digit. But let $U$ be the perfect unshuffle in base~$q$
on one index, which moves every digit to the next position and is a linear-time operation on tapes, and let $G_j$ be
the rotation of digit~$j$ of that index. Then $UG_j=G_{j+1}U$, so $[U,G_j]=(G_{j+1}-G_j)U$ has rank $d-d/q$, and
$\rank[1\otimes U,1\otimes G_j]=d^2-d^2/q$ for every~$j$ at once: one operation supplies the whole requirement for all
digits of that index. The requirement is real at every scale and can be met at all scales together.
\end{remark}

\subsection{What the theorem says}\label{sec:symmetry-says}

\begin{itemize}[leftmargin=1.5em]
\item \emph{For translation-covariant rules, the half in the half bound is a pigeonhole.} A code larger than half the
space meets its translate in at least $2D-d^2$ dimensions, and translation-covariant data has $d$ dimensions per uniform
piece. The argument does not depend on the characteristic, on the degree, on the window of the decoder, or on the other
tensor side. It does depend on exact covariance under a unit shift inside each class, on the classes being the same for
every output row, on linearity, and on there being one hop. The flag bounds of \cref{sec:order} are not shown to be
instances of it.
\item \emph{No symmetry that the swap lacks.} By \cref{lem:pigeonhole}, a code that exceeds the pairing value by
$\eps d^2$ needs $\rank[K,T]+\dim\ker[\sw,T]\ge d+2\eps d^2$ for every operator~$T$. Together with \cref{sec:order}:
to beat the pairing rate, the decoder may have neither a hidden order nor a symmetry, on either index, that the swap
does not share.
\item \emph{It uses the rule, not the relation.} The hypothesis of \cref{thm:TP} is that translated receivers share a
decoder. That is the feature of machines singled out in \cref{sec:rule}, and it is invisible in the cone relation.
\item \emph{The folklore step is optimal in dimension.} With $M$ equal pieces, nothing beats the $M\times M$ coarse
transposition by even one dimension, although many other codes equal it (\cref{rem:extremal}).
\end{itemize}

\section{The construction side}\label{sec:construction}

This section says what the two theorems exclude, what they leave, and what would have to be added to make them
statements about machines. Readings of the one-hop model as statements about tapes are labelled as readings.

\subsection{Mechanisms that fail}\label{sec:fail}

\Cref{tab:fail} lists upper bounds on the code dimension for decoders with structure on one index, in the linear
one-hop model, whatever is done on the other index. Each row is a theorem of \cref{sec:order,sec:symmetry}.

\begin{table}[ht]
\centering\small
\begin{tabular}{@{}>{\raggedright\arraybackslash}p{0.47\textwidth}>{\raggedright\arraybackslash}p{0.23\textwidth}l@{}}
\toprule
Structure on one index & Upper bound & By \\
\midrule
At most two Kronecker terms & $d(d+1)/2$ & \cref{cor:two-terms} \\
Commuting maps: any convolutions of any number of rows read in step & $d(d+1)/2$ & \cref{cor:commutative} \\
One $p$-group in characteristic~$p$: one dyadic hierarchy over $\F_2$ & $d(d+1)/2$ & \cref{cor:commutative} \\
All affine index maps, $d=p^m$, characteristic $p$ & $d(d+1)/2$ & \cref{cor:affine} \\
Expansion factor $c$ along some flag & $(1-\frac1{2c})d^2+O(d)$ & \cref{cor:expansion} \\
Uniform decoders with $M$ interval pieces, cycle or line & $(d^2+Md)/2$ & \cref{thm:TP,thm:line} \\
Copy rules with $M$ decoder classes of any shape & $(d^2+Md)/2$ & \cref{prop:copy} \\
$s$ modulation frequencies & $\frac s{s+1}(d^2+d)$ & \cref{cor:weyl} \\
One address digit of size $q$ treated uniformly & $\frac12(d^2+d^2/q)$ & \cref{cor:digit} \\
\bottomrule
\end{tabular}
\caption{Upper bounds on the code dimension in the linear one-hop model. The other index is arbitrary.}
\label{tab:fail}
\end{table}

Two design rules follow, for the one-hop model.
\begin{itemize}[leftmargin=1.5em]
\item \emph{Synchronous access is the pairing.} Reading any number of source rows in step with the output, at any
offsets and through any convolution, gives the symmetric matrices and no more.
\item \emph{Uniform passes are the pairing.} A stretch on which translated outputs share their decoder gains nothing,
whatever its window. The bound rises only with changes of regime, by $d/2$ symbols for each.
\end{itemize}

\subsection{A reading for tapes}\label{sec:reading}

The following is a reading, not a theorem about machines. A pass in which the position read and the position written
advance at the same speed, by a linear finite-state rule, computes a convolution with an unbounded window. On a line
this is a two-sided rule in the sense of \cref{thm:line}, with one class per stretch of unchanged regime, under two
provisos: the regime changes at the same positions for every output row (change points that depend on the row need
their common refinement), and the rule of each regime is applied to the whole stored word, with no state carried
across a change. A state of $q$ symbols carried across a change alters $K$ by a rank of at most $q$ per output row, so
by \cref{lem:pigeonhole} it costs a further term of order $qd$ per change. Under the provisos, \cref{thm:line} bounds
the gain over the pairing value by $d/2$ symbols per regime change, so rate $\tfrac12+\delta$ needs at least $2\delta d$
regime changes, and for equal pieces the bound is attained by the block-symmetric matrices, that is, by the coarse step
of the folklore recursion. Other speed ratios are decimations. On a cycle a single ratio is one twist class of
\cref{thm:twist}; mixtures of ratios fall under \cref{cor:expansion} or \cref{cor:affine}.

This reading matches what the series proved about motion (\cref{sec:motion}): locally affine motion is excluded on the
geometric side by the second and fourth papers, and here uniform decoding is excluded on the coding side, in a model
that has no geometry at all. The seventh paper has a statement of the same kind for storage codes on relations: Cayley
relations $\mathrm{Cay}(G,\{s,s^{-1},t\})$ without transitive triangles give no gain over the pairing rate, for every
finite group~\cite[Corollary~5.12]{SeventhPaper2026}. \Cref{thm:TP} allows any degree and any window, but only one hop and only linear
decoders.

\subsection{What a linear-time machine would need}\label{sec:would-need}

Combining \cref{sec:needs} with \cref{tab:fail}: in addition to positive characteristic, curved motion, expanding cone
relations and growing paired distance, a linear-time transposer, read in the one-hop model, needs on \emph{both}
indices a family of maps that
\begin{itemize}[leftmargin=1.5em]
\item has no invariant flag and no pair of flags with bounded expansion (\cref{thm:FL});
\item commutes with no non-derogatory translation up to $o(d)$ boundary columns, and intertwines none with boundedly
many pairwise commuting others (\cref{thm:TP,thm:twist});
\item treats no address digit uniformly (\cref{cor:digit}).
\end{itemize}
The digit masks at every scale, which the folklore recursion uses, have these properties: their number grows with
$d$, and no shift respects them. A bounded number of classes is not enough for rate near one, whatever their shape,
because $M$ classes leave a flag with $g(i)\le Mi$ and hence rate at most $1-1/(2M)+O(1/d)$. What the theorems leave
open is how large a code a given number of classes can carry when the classes have no almost-period. It can exceed the
value $(d^2+Md)/2$ of every copy rule and of every interval rule, at $d=6$ though not at $d=5$
(\cref{prop:aperiodic,rem:small}); whether a gain survives in rate is \cref{prob:aperiodic}.

\subsection{From the one-hop model to machines}\label{sec:to-machines}

Three steps separate \cref{thm:TP} from a statement about oblivious machines.
\begin{enumerate}[leftmargin=2em]
\item \emph{A hypothesis on cone systems.} The coding bounds of the series use properties of the cone relation, and
\cref{sec:false} shows that the relation-level properties proved so far cannot carry the conclusion. ``Translated receivers share their
decoder on every stretch of periodic head motion'' is a property of the rule. It has to be stated for the cone systems
of real schedules, where the receivers are output bits and the decoders are compositions of the transition function.
\item \emph{Depth reduction.} The genie of \cref{sec:depth} reveals values at chosen steps. The revealed word is advice
that depends on position, so it breaks uniformity wherever it is used. A version of \cref{thm:TP} with advice would have
to charge each revealed bit as a regime change; \cref{thm:TP} suggests the exchange rate, half a row per change, but
proves nothing about it.
\item \emph{Many hops.} \Cref{rem:compose-TP} bounds $t$ uniform hops by a loss that grows linearly in $t$. A machine
that runs for $cL$ steps has unboundedly many hops, and \cref{rem:compose-TP} shows that one cheap operation can break
every digit symmetry at once. So an argument for machines has to count something that an unshuffle does not supply for free.
\end{enumerate}
We do not know how to take any of these steps.

\section{Small sizes}\label{sec:small}

This section is numerical, except for \cref{prop:aperiodic}, which is a formal statement about one explicit code, and
for the results of \cref{sec:two-class} and of the subsequent memo~\cite{ClassUniformMemo2026} that are cited at the
end of \cref{rem:small} and at the end of the section. It reports exact computations over $\F_2$ for class-uniform rules
(\cref{def:class-uniform}) on a cycle. No proof in this paper uses them.

\subsection{Method}

For a partition of $\Z/d$ and a target dimension $D$, the existence of a class-uniform rule with a code of dimension at
least $D$ is a rank condition: the $d^2\times d^2$ matrix of $\sw-K$, whose row $(c,r)$ is determined by the functional
of the class of $c$ for the coordinate~$r$, must have rank at most $d^2-D$. We decide it in two ways.
\begin{itemize}[leftmargin=1.5em]
\item \emph{Exhaustively}, for $d\le4$, by enumerating subspaces and testing the criterion that a subspace $\X$ is a code
if and only if, for every class $J$, $\sum_{c\in J}\sigma^cX_c=0$ with $X_c\in\X$ implies
$\sum_{c\in J}(\text{row }c\text{ of }X_c)=0$ (\leanfile{tools/exhaustive_pigeonhole_checks.py}).
\item \emph{By a SAT solver}, for $d=5,6$: the rank condition is written as a factorization of the matrix through
$\F_2^{d^2-D}$, with the second factor in reduced echelon form, and given to
CryptoMiniSat (\leanfile{tools/class_uniform_sat.py}). A satisfying assignment is a rule together with its code; the
script rebuilds the matrix from the rule, recomputes its kernel, and \leanfile{tools/verify_class_uniform_code.py}
checks the saved code against the definition, sharing no code with the search. An answer ``unsatisfiable'' rests on the
solver, and we say so wherever we report one; as controls, the solver confirms the bound of \cref{thm:TP} for $d=4$ and
for $\{0,1\},\{2,3\},\{4\}$ at $d=5$. A time-out says nothing: the solver also times out on instances that are known to
be satisfiable. Some searches were restricted to codes that contain the space $\Hk^*$ of matrices with
$x_{c'}=\sigma^{c'-c}x_c$ for all $c,c'$ in one class, which is itself a code of dimension $Md$ and, for interval
classes, is the space $\Hk$ of \cref{sec:symmetry}. This makes the search faster and can only lower the result. The
code of \cref{prop:aperiodic} was found without the restriction and contains $\Hk^*$.
\end{itemize}

\subsection{Results}

\Cref{tab:small} lists the largest dimension of a code for partitions of $\Z/d$, up to the affine maps
$c\mapsto uc+t$ of the index. For $d\le5$ every partition with $1<M<d$ classes is listed; for $d=6$, the partitions
into two classes and the nine partitions into three classes for which the bound exceeds $(d^2+Md)/2$. ``Bound'' is the best value of \cref{thm:TP} over shift sets
of at most four elements, rounded down. Every code counted in the last column for $d\ge5$ is saved in
\leanfile{tools/data/}. The script \leanfile{tools/small_sizes.py} prints the table from these files, with the bounds
recomputed and every saved code checked again, and lists the partitions; \leanfile{tools/README.md} records the runs
of the solver, including those that timed out.

\begin{table}[ht]
\centering\small
\begin{tabular}{@{}cllccl@{}}
\toprule
$d$ & classes & pairing & $\lfloor(d^2+Md)/2\rfloor$ & bound & largest code \\
\midrule
3 & $\{0\},\{1,2\}$ & 6 & 7 & 7 & $=7$ \\
\midrule
4 & $\{0\},\{1,2,3\}$ & 10 & 12 & 12 & $=12$ \\
  & $\{0,1\},\{2,3\}$ & 10 & 12 & 12 & $=12$ \\
  & $\{0,2\},\{1,3\}$ & 10 & 12 & 12 & $=12$ \\
  & $\{0\},\{1\},\{2,3\}$ & 10 & 14 & 14 & $=14$ \\
  & $\{0\},\{2\},\{1,3\}$ & 10 & 14 & 14 & $=14$ \\
\midrule
5 & $\{0\},\{1,2,3,4\}$ & 15 & 17 & 17 & $=17$ \\
  & $\{0,1\},\{2,3,4\}$ & 15 & 17 & 17 & $=17$ \\
  & $\{0\},\{1,2\},\{3,4\}$ & 15 & 20 & 20 & $=20$ \\
  & $\{0\},\{1\},\{2,3,4\}$ & 15 & 20 & 20 & $19^{*}$ \\
  & $\{0\},\{1,4\},\{2,3\}$ & 15 & 20 & 21 & $20^{*}$ \\
  & $\{0\},\{1\},\{2\},\{3,4\}$ & 15 & 22 & 22 & $=22$ \\
\midrule
6 & $\{0\},\{1,\dots,5\}$ & 21 & 24 & 24 & $=24$ \\
  & $\{0,1\},\{2,\dots,5\}$ & 21 & 24 & 24 & $=24$ \\
  & $\{0,1,2\},\{3,4,5\}$ & 21 & 24 & 24 & $=24$ \\
  & $\{0,2,4\},\{1,3,5\}$ & 21 & 24 & 24 & $=24$ \\
  & $\{0,2\},\{1,3,4,5\}$ & 21 & 24 & 27 & $\ge24$ \\
  & $\{0,3\},\{1,2,4,5\}$ & 21 & 24 & 27 & $\ge23$ \\
  & $\{0,1,3\},\{2,4,5\}$ & 21 & 24 & 28 & $\ge25$ \\
  & $\{0,1\},\{2,4\},\{3,5\}$ & 21 & 27 & 30 & $\ge28$ \\
  & $\{0,1\},\{2,5\},\{3,4\}$ & 21 & 27 & 30 & $\ge28$ \\
  & $\{0,2\},\{1,4\},\{3,5\}$ & 21 & 27 & 30 & $\ge28$ \\
  & $\{0\},\{1,2,4\},\{3,5\}$ & 21 & 27 & 30 & $\ge28$ \\
  & $\{0\},\{1,5\},\{2,3,4\}$ & 21 & 27 & 30 & $\ge27$ \\
  & $\{0\},\{1,2,5\},\{3,4\}$ & 21 & 27 & 30 & $\ge27$ \\
  & $\{0\},\{1,3,4\},\{2,5\}$ & 21 & 27 & 30 & $\ge27$ \\
  & $\{0\},\{4\},\{1,2,3,5\}$ & 21 & 27 & 30 & $\ge27$ \\
  & $\{0\},\{3\},\{1,2,4,5\}$ & 21 & 27 & 30 & $\ge26$ \\
\bottomrule
\end{tabular}
\caption{Largest codes of class-uniform rules over $\F_2$. ``$=$'' marks an exact value: a code of that dimension
was found, and it equals the bound of \cref{thm:TP}. ``$\ge$'' marks the largest code found. A star marks a value that
was found and whose optimality rests on the solver (\cref{rem:small}).}\label{tab:small}
\end{table}

\begin{remark}[Numerical: what the table shows]\label{rem:small}
(a) \emph{Intervals.} For interval classes the bound $\lfloor(d^2+Md)/2\rfloor$ of \cref{thm:TP} is attained in every
case computed but one. For unequal pieces the codes that attain it are not block-symmetric, and at $d=4$ with pieces
$\{0\},\{1,2,3\}$ no copy rule attains it (\cref{rem:extremal}). The exception is $\{0\},\{1\},\{2,3,4\}$ at $d=5$: a
code of dimension $19$ exists; by \cref{lem:extremal} a code of dimension $20$ would contain $\Hk$; and the solver
reports that no code of dimension $20$ contains~$\Hk$. So the largest code has dimension $19$, on the word of the
solver, and the bound of \cref{thm:TP} is not attained for every partition into intervals.
(b) \emph{No almost-period, $d=5$.} For $\{0\},\{1,4\},\{2,3\}$ every single shift leaves four components, and the
bound of \cref{thm:TP} is $21$, from a set of three shifts. A code of dimension $20=(d^2+Md)/2$ exists, and the solver reports that
none of dimension $21$ does.
(c) \emph{No almost-period, $d=6$.} For $\{0,1,3\},\{2,4,5\}$ every nonzero shift leaves four components, and the best
bound of \cref{thm:TP} is $28$, from a set of three shifts. A code of dimension $25$ exists, one more than
$(d^2+Md)/2=24$
(\cref{prop:aperiodic}); searches for dimension $26$ timed out. Four partitions into three classes have codes of
dimension $28$, one more than $(d^2+Md)/2=27$.
(d) The gain over $(d^2+Md)/2$ in (c) is one dimension in each case. The data, at $d\le6$, do not distinguish a gain
of bounded size from one that grows with~$d$. For every fixed number $M\ge2$ of classes it can grow, and this is proved, not
numerical, though for other partitions than those of the table. For two classes the gain is at least $d/4-1$ at every
$d=4k\ge8$, over every field (\cref{prop:two-class}, which extends \cite[Theorem~4.8]{ClassUniformMemo2026} to
characteristic~$2$); the codes over $\F_2$ at $d=8$, $12$ and $16$ are saved with those of \cref{tab:small}
(\cref{rem:two-class}). For every $M\ge2$ there are partitions into $M$ classes with a gain of order $Md/8$
(\cref{cor:every-M}). And when the number of classes grows in proportion to $d$, the blow-ups of the code of
\cref{prop:aperiodic} are codes of dimension $25k^2$ for $2k$ classes on $\Z/6k$, a gain
of~$k^2$~\cite[Theorem~4.2 and Corollary~4.4]{ClassUniformMemo2026}.
\end{remark}

\begin{proposition}[Two classes with no almost-period beat every interval rule]\label{prop:aperiodic}
\formal{AperiodicExample.lean}{exists_classUniformCode_gt,isClassUniformCode_aperiodicCode,finrank_aperiodicCode,aperiodicCode,aperiodicCls_eq_true_iff}
For $d=6$ over $\F_2$ and the classes $\{0,1,3\}$ and $\{2,4,5\}$ of $\Z/6$ there is a code of a class-uniform rule of
dimension $25$. In particular the bound $(d^2+Md)/2$ of \cref{thm:TP} for $M$ interval classes, and of \cref{prop:copy}
for copy rules with $M$ classes of any shape, does not hold for coded rules on arbitrary classes.
\end{proposition}

\begin{outline}
The code is the span of $25$ explicit matrices, and the two decoders are explicit
(\leanfile{AperiodicExampleData.lean}, generated from a certificate by \leanfile{tools/gen_aperiodic_example.py}). That
every generator satisfies the rule at every row, $900$ equations over $\F_2$, is checked by the kernel
(\lean{aperiodicGen_row}); by linearity the span is a code (\lean{isClassUniformCode_aperiodicCode}). Each generator has
an entry equal to~$1$ where all the others vanish (\lean{aperiodicGen_pivot}), so the generators are independent and
the dimension is~$25$ (\lean{finrank_aperiodicCode}). Then $2\cdot25=50>48=6^2+2\cdot6$
(\lean{exists_classUniformCode_gt}). The certificate is used only as data.
\end{outline}

The code of \cref{prop:aperiodic} was found by the solver. It contains the $12$-dimensional space $\Hk^*$, it meets its translate by one column in exactly $2\cdot25-36=14$
dimensions, the least that \cref{lem:pigeonhole} allows, and it meets its translates by two and three columns in $18$
and $17$ dimensions (\leanfile{tools/analyse_code.py}). We do not have a description of the code itself, or a way to
enlarge it at a fixed number of classes, and one example at $d=6$ says nothing about
rates: the term $Md/2$ is a quarter of $(d^2+Md)/2$ at this size. What the example does show is that the union bound in
\cref{thm:TP} is not an artefact, and that masks without an almost-period carry a coding gain that no copy rule and no
interval rule has. \Cref{prob:aperiodic} asks how large it can be.

The subsequent memo~\cite{ClassUniformMemo2026} proves two further facts about this example, in Lean, and it finds
larger gains of a different kind.
\begin{itemize}[leftmargin=1.5em]
\item \emph{The gain does not come from characteristic~$2$.} The same two classes carry a code of dimension $25$ over
every field in which $3\ne0$, in particular over $\Q$~\cite[Theorem~3.2]{ClassUniformMemo2026}. Whether they do in
characteristic~$3$ is open.
\item \emph{The example can be blown up.} If a colouring of $\Z/d$ with $M$ classes carries a code of dimension $D$,
then a colouring of $\Z/dk$ with $kM$ classes carries a code of dimension
$k^2D$~\cite[Theorem~4.2]{ClassUniformMemo2026}. Blowing up multiplies the gain over $(d^2+Md)/2$ and the product $dM$
by the same factor $k^2$, so blowing up this one example gives growing gains only together with a growing number of
classes. Blowing up a family by a fixed factor is another matter (\cref{cor:every-M}).
\item \emph{Two classes can gain more than one dimension.} For $d=4k\ge8$ the classes $\{0,2,\dots,d-4\}$ and the rest
carry an explicit rule, each of whose decoders reads one anti-diagonal, with a code of dimension
$(d^2+2d)/2+d/4-1$. The memo proves this over fields in which $2\ne0$, in its text, and proves in Lean at $d=8$, $12$
and~$16$ that the classes carry a code of this dimension~\cite[Theorems~4.8 and~4.9]{ClassUniformMemo2026};
\cref{prop:two-class} proves it over every field, and \cref{cor:every-M} derives from it an unbounded gain for every
fixed number of classes. Unlike the code of \cref{prop:aperiodic}, these codes have a description, and it is one that
can be analysed.
\end{itemize}

\section{Open problems}\label{sec:open}

The question~(T) is open. For oblivious machines it is the question about the function $r(\eps)$ of
\cref{sec:one-function}. The problems below are the precise points at which this paper stops; the first three are finite
or algebraic and can be attacked directly.

\begin{problem}[Classes with no almost-period]\label{prob:aperiodic}
For a partition $\mathcal P$ of $\Z/d$ into $M$ classes let $D(\mathcal P)$ be the largest dimension of a code of a
class-uniform rule (\cref{def:class-uniform}), over a given field. \Cref{thm:TP} gives $D(\mathcal P)\le(d^2+Md)/2$ when
the classes are intervals. For general classes this fails (\cref{prop:aperiodic}), and for every fixed number $M\ge2$
of classes it fails by an amount that grows with $d$ (\cref{prop:two-class,cor:every-M}).
\begin{enumerate}[label=(\alph*),leftmargin=2em]
\item For two classes, is $D(\mathcal P)\le(\tfrac12+o(1))d^2$? More generally, is this so for $o(d)$ classes?
For two classes of any shape the best upper bound known is
$\tfrac{41}{64}d^2+\tfrac53d^{3/2}$~\cite[Theorem~3.4]{SharedDecoders2026}; before that paper it was
$\tfrac34d^2+O(d)$, from \cref{thm:FL}, together with what \cref{thm:TP} gives for the best shift set. For $M\le d/2$
classes that paper proves the bound
$\bigl(1-1/(4\lceil\sqrt{2M}\,\rceil)\bigr)d^2$~\cite[Theorem~3.5]{SharedDecoders2026}, which is below the value
$(1-\tfrac1{2M})d^2+O(d)$ of \cref{thm:FL} when $M\ge11$. The largest gains known for a fixed number of
classes are of order~$d$.
\item Determine $D$ for the classes $\{0,1,3\},\{2,4,5\}$ of $\Z/6$. By \cref{thm:TP} it is at most $28$ over every
field. It is at least $25$ over every field in which $3\ne0$: over $\F_2$ by \cref{prop:aperiodic}, and in general
by~\cite[Theorem~3.2]{ClassUniformMemo2026}. In characteristic~$3$ no code of dimension $25$ is known.
\item Fix the number $M\ge2$ of classes. Is there a constant $C_M$ such that $D(\mathcal P)\le(d^2+Md)/2+C_Md$ for
every $d$ and every partition $\mathcal P$ of $\Z/d$ into $M$ classes? Or is there a family of partitions into $M$
classes, with codes that have a description that can be analysed, in which $D(\mathcal P)-(d^2+Md)/2$ grows faster
than~$d$?
\end{enumerate}
\end{problem}

A positive answer to (a) would be the half bound for all piecewise-uniform decoders in the one-hop model; it needs a
count sharper than the union bound of \cref{thm:TP}, which loses a factor $s/(s+1)$ for $s+1$ translates. A negative
answer would be a coding gain of positive rate over the pairing value from $o(d)$ classes, native to masks and
rotations, the first of its kind, and its cost on a tape would become the question. \Cref{sec:small} records what exact
computation says for small~$d$.

The two bounds quoted in~(a) come from a later paper of the series~\cite{SharedDecoders2026}, which proves them in
Lean, over every field, for codes that are subspaces and decoders that are arbitrary maps (\cref{sec:sets}). Its count
is sharper than the union bound in the way asked for above. The code is intersected with one translate, the
intersection lies in the space of matrices that are covariant on the components of that shift, this space is
invariant under rotation, and the next intersection is taken inside it; so with shift sets $E_0=\{0\}$ and
$E_j=E_{j-1}\cup(E_{j-1}+t_j)$ for $1\le j\le k$, a code has dimension at most
$\sum_{j<k}n_j/2^{j+1}+n_k/2^k$, where $n_j$ is $d$ times the number of components for
$E_j$~\cite[Theorem~3.2]{SharedDecoders2026}. When $E_2$ has four elements, \cref{thm:TP} for $E_2$ gives
$\tfrac34n_0+\tfrac14n_2$ where this gives $\tfrac12n_0+\tfrac14n_1+\tfrac14n_2$, and $n_0=d^2$. The value $\tfrac{41}{64}$ comes from the shifts $s,3s,9s$ and an average over~$s$. For
almost every partition into two classes the same paper has a slightly smaller constant, and it proves that the
certificates of this method cannot go below $(\tfrac58-o(1))d^2$, and those of a variant not below
$(\tfrac9{16}-o(1))d^2$, for almost every such partition~\cite[Theorems~C.5, C.6 and~C.7]{SharedDecoders2026}. So (a) is open, and going below
these values needs more than the certificates of that method.

Part~(c) has a history. Revision \texttt{b2c389a} of this paper asked only for a family of partitions and codes, with
a description that can be analysed, in which $D(\mathcal P)-(d^2+Md)/2$ grows, without fixing the number of classes.
The memo~\cite{ClassUniformMemo2026} answered that question as it was worded: the blow-ups of the code of
\cref{prop:aperiodic} are codes of dimension $25k^2$ for $2k$ classes on $\Z/6k$, so that
$D(\mathcal P)-(d^2+Md)/2\ge k^2$ there~\cite[Theorem~4.2 and Corollary~4.4]{ClassUniformMemo2026}. The family is
described given its base code, which a solver found. A blow-up multiplies the excess of a code over $(d^2+Md)/2$ and
the product $dM$ by the same factor~\cite[Corollary~4.3]{ClassUniformMemo2026}, so the blow-ups of one finite example
gain only together with a growing number of classes. The memo and revision \texttt{14b157a} of this paper therefore
asked whether $D(\mathcal P)-(d^2+Md)/2$ is bounded for a fixed number of
classes~\cite[Problem~6.2]{ClassUniformMemo2026}. It is not. An update of the memo showed this for $M=2$ over every
field in which $2\ne0$, with an explicit rule whose code has a gain of
$d/4-1$~\cite[Theorem~4.8 and Corollary~4.10]{ClassUniformMemo2026}, and left open characteristic~$2$ and every
$M\ge3$. \Cref{prop:two-class} is that theorem with its two coefficients $2$ and $\tfrac12$ replaced by $\lambda$ and
$\lambda^{-1}$, which makes it hold over every field. \Cref{cor:every-M} blows the two-class codes up by a fixed
factor, with a code of its own at each position, and gives an unbounded gain for every $M\ge2$ over every field. What
remains is the rate of growth, which is~(c) as now stated; for $M=2$ it is, in substance,
\cite[Problem~6.3]{ClassUniformMemo2026}.

By \cref{cor:every-M} a constant $C_M$ as in~(c) is at least $\lfloor M/2\rfloor/4$, over every field; for two classes
it is at least $\tfrac14$. For the two classes of \cref{prop:two-class}, \cref{thm:TP} bounds the gain by $d/2$. In
general \cref{thm:TP} does not answer~(c): with two shifts it bounds the gain of a two-class partition by
$(\chi-2)d/2$, where the number $\chi$ of components can be of order~$d$, and with three or more shifts its bound on
$D(\mathcal P)$ is at least $2d^2/3$~\cite[Section~6]{ClassUniformMemo2026}. A positive answer to~(c) for $M=2$ would
prove (a) for two classes. A family of two-class partitions with a gain of order $d^2$ would be a negative answer
to~(a).

For partitions with a bounded period the answer to~(c) is positive, and for some of them the constant is known. Let a
partition $p$ of $\Z/s$ be lifted to the partition $p^\uparrow$ of $\Z/d$, $d=ms$, in which $c$ has the class of
$c\bmod s$. \Cref{thm:TP} with the shifts $0$ and $s$ bounds its largest code by $(d^2+sd)/2$. If $m$ is prime to $s$
and not divisible by the characteristic, and the field contains the $m$-th roots of unity, the exact value is
\[
  D(p^\uparrow)=\frac{d^2-ms^2}2+m\,D(p)=\frac{d^2}2+\frac ds\Bigl(D(p)-\frac{s^2}2\Bigr);
\]
over a field without these roots of unity the right side is a lower bound, and it is an upper bound when $D(p)$ is
taken over the algebraic closure~\cite[Theorems~5.1 and~5.3]{SharedDecoders2026}. So the gain of a partition with
period $s$ is at most $(s-M)d/2$, for these lifts it is a fixed multiple of~$d$, which is zero when $D(p)$ is the
value $(s^2+Ms)/2$ of intervals, and a gain that grows faster than $d$ needs partitions without a bounded period; the
classes of \cref{prop:two-class} have no period smaller than~$d$. From the classes of \cref{prop:aperiodic} the lift
gives two classes on $\Z/6m$, for $m$ prime to~$6$, with a binary code of dimension $d^2/2+7d/6$, and this is the
whole code of the lifted rule~\cite[Corollary~5.2 and Theorem~B.2]{SharedDecoders2026}. Its gain over $(d^2+2d)/2$
is $d/6$; the rule of \cref{prop:two-class} gains $d/4-1$, at the sizes divisible by four.

\begin{problem}[Unequal pieces]\label{prob:unequal}
For which partitions of $\Z/d$ into $M$ intervals, and over which fields, is the bound $\lfloor(d^2+Md)/2\rfloor$ of
\cref{thm:TP} attained? It is attained for equal pieces (\cref{prop:block-symmetric}). Over $\F_2$ it is attained in
every case computed in \cref{sec:small} except one, the pieces $\{0\},\{1\},\{2,3,4\}$ at $d=5$, where the largest code
has dimension $19<20$ on the word of a solver. Some of the codes that attain it are not codes of copy rules.
\end{problem}

The question is about linear decoders. If the decoders may be arbitrary maps and the code is still a subspace, the
bound is attained over $\F_2$ for every partition into $M$ intervals with $d-M$ even, and the same construction gives
a code of dimension $(d^2+Md)/2$ for every partition into $M$ classes of any shape with $d-M$
even~\cite[Theorem~4.3]{SharedDecoders2026}. For the pieces $\{0\},\{1\},\{2,3,4\}$ at $d=5$ this is a code of
dimension $20$ with decoders that are not linear, where linear decoders reach $19$ on the word of a solver.

\begin{problem}[Both sides, all scales]\label{prob:both}
Let both the $A_a$ and the $B_a$ lie in the span of masks of dyadic intervals composed with rotations, with at most $M_j$
intervals of length $d/2^j$ in use at scale~$j$. Bound $\dim\ker(\sw-K)$ in terms of the profile $(M_j)_j$. The folklore
recursion shows that the full profile suffices for dimension $d^2$; \cref{thm:TP} and \cref{cor:digit} treat one side
and one scale at a time.
\end{problem}

\begin{problem}[Shared decoders on cone systems]\label{prob:machines}
\begin{enumerate}[label=(\alph*),leftmargin=2em]
\item State and prove a version of \cref{thm:TP} for the cone systems of oblivious linear-time schedules: on a stretch
of periodic head motion the outputs share their decoder, and a revealed bit or an irregular step changes the decoders
of some outputs. Does a bound of the form $r(\eps)\le\tfrac12+C(\eps+\eta)$ hold for schedules whose motion is
periodic outside a set of steps of density~$\eta$? (\Cref{sec:to-machines}.)
\item Fix a window $w\ge3$ and an alphabet of $q$ symbols. Does every storage code on a one-class relation of
window~$w$ (\cref{sec:sets}) have at most $q^{(1/2+o(1))d^2}$ elements as $d$ grows, or even at most $q^{d^2/2+Cd}$
with a constant $C$ that depends on $w$ and~$q$? And is this so at least when one function serves all rows, that is,
when the function for the position $(c,r)$ does not depend on~$c$? (For $w\le2$ the answer to both is yes, with
$C=\tfrac32$: \cref{prop:order}(b). For $w=2$ and a function for each position the least constant lies between $1$
and $\tfrac32$, by the examples after that proposition; for $w=1$ it is $\tfrac12$.)
\item Fix $w\ge3$ and an alphabet of $q$ symbols, and let $\rho$ assign to every column index at most $w$ rows. Does
every $\rho$-code (\cref{sec:sets}), that is, every set of matrices in which each column is a function of its $w$
rows, whole, have at most $q^{(1/2+o(1))d^2}$ elements as $d$ grows? (For $w\le2$ it has at most
$q^{d(d+1)/2+d-1}$, by the proof of \cref{prop:order}(b); in general at most $q^{wd(d+1)/(w+1)}$, by
\cref{prop:rows}(d).)
\end{enumerate}
\end{problem}

Part~(a) is the problem as the earlier revisions stated it; part~(b) was added in revision \texttt{3911d87}, for
every window, and is restricted here to the windows that remain open. The form of the bound
in~(a) comes from the one-hop model with affine rules. By \cref{cor:entry-advice}, over a finite field, if every rule
that the advice selects changes, from one row to the next, at no more than $E$ entries, then the matrices of a code
that share a value of the advice number at most $\abs F^{(d^2+d+E)/2}$, whatever the advice is. The rate of such a
set, the analogue of $r(\eps)$, is therefore at most $\tfrac12+(d+E)/(2d^2)$, and $E=O((\eps+\eta)d^2)$ gives a bound
of the form above. So the question is whether the revealed bits and the irregular steps of a schedule amount to that
many changes of decoders at single outputs, on the cone systems of real schedules and for decoders that are
compositions of a transition function. Revision \texttt{b2c389a} of this paper stated the count as ``a genie that
reveals $\eps L$ bits changes the regime at most $\eps L$ times'', and its \cref{sec:to-machines} priced a change of
regime at half a row. Half a row is the price of a change that affects all $d$ entries of a row, and at that price the
count gives nothing.

For decoders that are not affine the count of \cref{cor:entry-advice} is not available at all, and this is the
reason for part~(b). With a view of the whole matrix, one decoder for all rows, with no change of any kind, is correct
on a set of at least $q^{d^2}/(2d)$ matrices~\cite[Theorem~A.1 and Corollary~A.2]{SharedDecoders2026}
(\cref{sec:sets}). The inputs that share a value of advice form a set and not a subspace. So for them, and for
decoders that are compositions of a transition function, the sharing of a decoder bounds nothing when the decoder
sees the whole stored word. What that example uses, and the decoders of a machine do not have, is the size of the
view. The cones that a genie leaves have bounded size, so a decoder reads boundedly many bits, and with them it can
tell apart only boundedly many positions (\cref{sec:sets}). Whether the sharing of a decoder bounds a set of matrices
when the decoders read boundedly many symbols of a fixed alphabet is not known, and part~(b) asks it in the one-hop
model, in two forms. In the weaker form one function serves all rows and reads translated positions, which is what a
stretch of periodic motion gives. Since a set of matrices can contain marker rows, one function can then imitate
boundedly many classes of rows, where the marker rows are free: with markers the value $(d^2+d)/2$ of subspaces is
exceeded by a multiple of $d$ that grows with the window, over every finite field (\cref{sec:sets}), so the constant
$C$ of~(b) has to grow with~$w$. In the stronger form only
the positions are translates, and every position may have its own function. Over a fixed alphabet there are boundedly many functions of $w$ symbols, so the stronger form
is the question for boundedly many classes of any shape, drawn entry by entry. It contains what is left of the weaker
form among the inputs that share a revealed word: there the outputs of a periodic stretch whose cones meet the same
pattern of revealed bits, with the same values, compute the same function of the input bits that they read, so that
the revealed word divides the outputs into boundedly many classes of a shape that the advice chooses. The hypothesis
of the stronger form is a property of the relation alone, and so of the head motion. What is proved is in
\cref{prop:sets,prop:order,prop:mutual}: windows one and two; every window when the offsets, or the differences of
the rows from a permutation, are bounded or take boundedly many values; codes that use only helpers in both
directions; and, for the weaker form, subspaces. What these
results leave is a window of at least three on a relation whose rows and offsets are both spread. For small windows
the acyclic-subgraph bound, which gives all of them but the last, may still be enough; for large windows the computations
reported in \cref{sec:sets}, in the paragraph on what the order argument does not reach, indicate that it is not. A code above the pairing rate would have to use
helpers in one direction, on a relation of that kind.

The conclusion asked for in~(b) is that of the all-code question of the fifth
paper~\cite[Question~9.1]{DirectedRecovery2026}, with other hypotheses. That question is about relations of bounded
incidence degree that are separated, and for those the answer is no (\cref{sec:false}). A one-class relation has
the symmetry of its positions instead; it may have loops and cycles of length two, and one position may help $d$
others. We have not determined whether relations like those of the counterexamples of \cref{sec:false} can be
one-class relations; the undirected ones cannot, in the sense of \cref{prop:mutual}(b). In every one-class
relation of bounded window few positions are their own helpers, and $o(d^2)$ positions meet all transitive triangles
(\cref{prop:mutual}(c)). A rule with $M$ classes of any shape whose decoders have window $w$ gives a relation that is
contained in a one-class relation of window $Mw$. So a positive answer to the stronger form of~(b) would give the
half bound of \cref{prob:aperiodic}(a) for every fixed number of classes and decoders of bounded window over a fixed
finite field, for sets as well as subspaces. For subspaces the weaker form is enough. A code of a rule with $M$
classes whose decoders have window $w$ can be replaced by a code of at least the same dimension in which one row of
each class is free, with decoders of window at most
$2w+1$~\cite[Theorem~4.1]{SharedDecoders2026}; fixing $k=\lceil\log_qM\rceil$ of these rows as markers
(\cref{sec:sets}) gives a storage code with one function for all rows, on a one-class relation of window
$M(2w+1)+k$, with $q^{-kd}$ times as many elements. So a positive answer to the weaker form, for every window, gives
\cref{prob:aperiodic}(a) for rules with a fixed number of classes and decoders of bounded window over a fixed finite
field. A negative answer to either form would be a code above the pairing rate, with bounded windows, on positions
that move with the output row.

Part~(c) was added with \cref{prop:rows,prop:stages,prop:expanding,prop:pieces}. A storage code on a one-class
relation of window $w$ is a $\rho$-code with at most $w$ rows for each column, whatever its offsets, so a positive
answer to~(c) for every $w$ gives the bound $q^{(1/2+o(1))d^2}$ of~(b) for every window; it does not give the form
$d^2/2+Cd$. In the other direction (c) allows more than (b): an entry may depend on all entries of its rows. What is
proved is the case in which the rows of each column lie in a small block, and with it pipelines of stages that read
rows inside address cosets, where $t$ moved digits give exactly $(d^2+2^td)/2$; pipelines with buffered scans and
with differently arranged copies, more weakly (\cref{prop:stages}); and products of codes on boundedly many positions
(\cref{prop:pieces}). What is left are the relations on which the acyclic-subgraph bound stays a constant above one
half whatever order is used, among them those that expand and contain a bijection. There are such relations that are
unions of the identity and three increasing maps, so that an operator with these rows takes boundedly many forward
passes, given the strings that mark the maps (\cref{prop:expanding} and the paragraph after it). A positive answer
to~(c) for $w\ge4$ therefore needs an argument that is
not an order argument, and a negative answer on increasing maps, with functions computable in linear time, would be
an operator of boundedly many passes that transposes more than the pairing rate. By \cref{prop:pieces} a code above
the pairing rate cannot be a product of local codes. One family in which to look for either answer: for a group
$\Gamma$ of order $d$ and a bounded set $\Sigma\subseteq\Gamma$, index the rows and columns by $\Gamma$ and let the
helpers of $(c,r)$ be the positions $(rs,ct)$ for boundedly many pairs $(s,t)$ with $s\in\Sigma$; these relations
are contained in the relation of $\rho(r)=r\Sigma$, and for $\Gamma=\Z/d$ they are the one-class relations with
constant offsets and constant row differences, which \cref{prop:order}(d) settles.

For adaptive machines a later paper of the series puts the function $r(\eps)$ on the machine's own runs. At the
canonical entry width, for all but an exponentially small share of the inputs that a linear-time transposer handles,
and for every set $B$ of steps of the run on such an input, the relation between entries that is left when $B$ is cut
out carries storage codes whose deficit is at most a small multiple of the input length plus a constant times
$\card B$~\cite[Theorem~H.5 and Corollary~H.6]{AdaptiveTransfer2026}. So a bound on storage codes that is stated for
relations, as part~(b) is, can be applied to a single run of an adaptive machine wherever the relation of that run
satisfies its hypothesis.

\begin{problem}[What an unshuffle does not supply]\label{prob:scales}
\Cref{cor:digit} requires a transposer to break the symmetry of every address digit, and \cref{rem:compose-TP} shows
that commutator rank cannot turn this into a factor $\log d$, because one unshuffle breaks all digits at once. Find a
quantity that is subadditive over the steps of a machine, is large for the swap at every scale separately, and is not
raised at all scales by a single linear-time operation. This is the operator form of the multi-scale problem of the
sixteenth paper~\cite[Problem~11.1]{SixteenthPaper2026}.
\end{problem}

For pipelines whose stages read rows inside cosets of the address space, the dimension of the coset adds up over
the stages and bounds the set that is transposed, for all codes: one for each digit moved, and $\log_2d$ for the
swap (\cref{prop:stages}(a)). This is not the quantity asked for. It does not look at scales separately, and a
single linear-time operation can raise it to its maximum: a scan with a buffer of one row, or a stage that reads a
row $r$ together with its shuffle parents $\lfloor r/2\rfloor$ and $\lfloor r/2\rfloor+d/2$ (the remarks after
\cref{prop:stages}). For the
first the proposition has a price in advice, and for the second only the weaker count by the number of ones of the
address.

\begin{problem}[Relative rigidity on ball patterns]\label{prob:rigidity}
The eighth paper asks whether there is $\rho>0$ such that, for every linear-time schedule, every fixed $D$ and all
large~$d$, the transposition matrix is not the sum of a matrix of rank below $\rho L$ and a matrix supported on the
depth-$D$ ball pattern of the schedule~\cite[Problem~10.6]{EighthPaper2026}; a positive answer alone would exclude
oblivious linear-time transposers with linear data paths~\cite[Corollary~9.3]{BoundedStack2026}. The classes of this
paper are defined by structure on one tensor side. Which ball patterns of schedules are contained in such a class,
after a change of bases on one side? For those, \cref{thm:FL,thm:TP} give $\rho=\tfrac12-o(1)$.
\end{problem}


\appendix
\section{Index of Lean declarations}\label{app:leanindex}

Every declaration of this paper's library cited in the text, grouped by the Lean module that contains it, relative to the namespace
\texttt{SwapRigidity}. All of them depend only on the axioms \texttt{propext}, \texttt{Classical.choice} and
\texttt{Quot.sound} (\cref{app:verification}).

\paragraph{\texttt{SwapRigidity.AperiodicExample}}
{\small\raggedright \texttt{aperiodicCls\_eq\_true\_iff}, \texttt{aperiodicCode}, \texttt{aperiodicGen\_pivot}, \texttt{aperiodicGen\_row}, \texttt{exists\_classUniformCode\_gt}, \texttt{finrank\_aperiodicCode}, \texttt{isClassUniformCode\_aperiodicCode}.\par}

\paragraph{\texttt{SwapRigidity.BlockSymmetric}}
{\small\raggedright \texttt{blockSymmetric}, \texttt{blockSymmetric\_attains\_bound}, \texttt{isClassUniformCode\_blockSymmetric}, \texttt{two\_mul\_finrank\_blockSymmetric}.\par}

\paragraph{\texttt{SwapRigidity.ClassUniform}}
{\small\raggedright \texttt{IsClassUniformCode}, \texttt{IsClassUniformCode.finrank\_le}, \texttt{IsClassUniformCode.row\_eq\_of\_rot\_mem}, \texttt{IsClassUniformCode.two\_mul\_finrank\_le}, \texttt{IsClassUniformCode.two\_mul\_finrank\_le\_of\_const}, \texttt{Link}, \texttt{rot}.\par}

\paragraph{\texttt{SwapRigidity.Flag}}
{\small\raggedright \texttt{card\_staircase}, \texttt{eq\_zero\_of\_mem\_swapCode\_of\_staircase}, \texttt{exists\_staircase\_finrank\_eq}, \texttt{finrank\_swapCode\_conj}, \texttt{finrank\_swapCode\_le\_of\_flag}, \texttt{finrank\_swapCode\_le\_of\_flag\_left}, \texttt{finrank\_swapCode\_le\_of\_staircase}, \texttt{finrank\_swapCode\_staircase}, \texttt{finrank\_swapCode\_swap}, \texttt{mem\_swapCode}, \texttt{mem\_swapCode\_staircase\_iff}, \texttt{staircaseRestrict\_injective}, \texttt{swapCode}, \texttt{swapCodeEquiv}, \texttt{two\_mul\_finrank\_swapCode\_le\_of\_triangular}, \texttt{two\_mul\_finrank\_swapCode\_le\_of\_triangular\_left}, \texttt{two\_mul\_sum\_fin\_succ}.\par}

\paragraph{\texttt{SwapRigidity.Line}}
{\small\raggedright \texttt{IsLineUniformCode}, \texttt{IsLineUniformCode.row\_succ\_eq}, \texttt{IsLineUniformCode.two\_mul\_finrank\_le}, \texttt{isLineUniformCode\_of\_coeff}, \texttt{isLineUniformCode\_symmetric}, \texttt{lineView}, \texttt{truncShift}.\par}

\paragraph{\texttt{SwapRigidity.Pigeonhole}}
{\small\raggedright \texttt{card\_mul\_finrank\_add\_le\_iInf\_comap}, \texttt{finrank\_ker\_commutator\_sub\_le}, \texttt{finrank\_range\_commutator\_comp\_le}, \texttt{finrank\_range\_commutator\_le}, \texttt{mul\_finrank\_le\_add\_finrank\_iInf\_comap}, \texttt{two\_mul\_finrank\_ker\_sub\_le}, \texttt{two\_mul\_finrank\_ker\_sub\_le\_ker}, \texttt{two\_mul\_finrank\_ker\_sub\_le\_of\_commute}.\par}

\paragraph{\texttt{SwapRigidity.Swap}}
{\small\raggedright \texttt{card\_le\_finrank\_ker\_swap\_commutator}, \texttt{two\_mul\_finrank\_ker\_swap\_sub\_one}, \texttt{two\_mul\_finrank\_range\_swap\_sub\_one}.\par}

\paragraph{\texttt{SwapRigidity.Twist}}
{\small\raggedright \texttt{finrank\_le\_sum\_finrank\_twistSpace}, \texttt{succ\_mul\_finrank\_swapCode\_le}, \texttt{succ\_mul\_finrank\_swapCode\_le\_of\_twist\_classes}, \texttt{twistSpace}, \texttt{two\_mul\_finrank\_swapCode\_le\_of\_twist}.\par}

\section{Building and verifying the formalization}\label{app:verification}

The Lean project is the directory \leanfile{lean/} next to this paper. It uses Lean and Mathlib, both at version v4.34.1
(\leanfile{lean-toolchain}, \leanfile{lakefile.toml}, \leanfile{lake-manifest.json}), and requires the shared
\leanfile{lean-machine/} library of the repository through a relative path, which supplies Mathlib at the pinned
revision. It uses no library of an earlier paper. Keep the repository's directory layout so that
\leanfile{lean-machine/} stays in place.

\paragraph{Build.} With \texttt{elan} installed, run in \leanfile{lean/}:
\begin{center}
\texttt{lake exe cache get \&\& lake build}
\end{center}
The first command downloads the compiled Mathlib cache; the second compiles \leanfile{SwapRigidity}.

\paragraph{Axioms and index.} The file \leanfile{lean/PaperAxioms.lean}, generated from the \LaTeX{} sources of this paper
by \leanfile{tools/lean_index.py}, runs \texttt{\#print axioms} on every declaration cited in this paper
(\cref{app:leanindex}):
\begin{center}
\texttt{lake env lean PaperAxioms.lean}
\end{center}
Every declaration depends at most on \texttt{propext}, \texttt{Classical.choice} and \texttt{Quot.sound}. No Lean file
of this library uses \texttt{sorry}, \texttt{admit} or \texttt{native\_decide}, declares an \texttt{axiom}, bypasses the
kernel via \texttt{unsafe}, \texttt{implemented\_by} or \texttt{extern}, or raises \texttt{maxHeartbeats}. Running
\texttt{python3 tools/lean\_index.py} from the paper's root directory regenerates both
\leanfile{paper/sections/leanindex.tex} and \leanfile{lean/PaperAxioms.lean} (\texttt{--index-only} writes only the
index; \texttt{--check} verifies that every cited name resolves and that both generated files are current).

\paragraph{Where things are.} The library \leanfile{lean/SwapRigidity/} has the following modules, all imported by the
root module \leanfile{lean/SwapRigidity.lean}.
\begin{itemize}[leftmargin=1.5em]\raggedright
\item \leanfile{Flag.lean} (\cref{sec:order}): the swap code, the elimination on staircase patterns, invariance under
changes of bases, the coordinate form of Theorem~FL, the corollary for a pair of flags, the exchange of sides, and
attainment.
\item \leanfile{Pigeonhole.lean} (\cref{sec:symmetry}): the pigeonhole lemma in its three forms, the rank bound for
commutators and its subadditivity under composition, and the dimension count for an intersection of preimages.
\item \leanfile{Swap.lean}: the dimensions of the kernel and image of $\sw-1$ and the bound $\dim\ker[\sw,T]\ge d$.
\item \leanfile{ClassUniform.lean}: class-uniform codes and Theorem~TP on a cycle.
\item \leanfile{Entry.lean} (\cref{prop:entry}): rules given entry by entry, and the bound for decoders that change at
single entries.
\item \leanfile{EntryAdvice.lean} (\cref{cor:entry-advice}): the same rules with advice of any kind, over a finite
field.
\item \leanfile{BlockSymmetric.lean}: the block-symmetric code, its dimension, and attainment of the bound.
\item \leanfile{Line.lean}: two-sided rules on a line and the line form of Theorem~TP.
\item \leanfile{Twist.lean}: the twist-class theorem.
\item \leanfile{AperiodicExample.lean} and the generated \leanfile{AperiodicExampleData.lean} (\cref{prop:aperiodic}):
an explicit class-uniform code of dimension $25$ for $d=6$ over $\F_2$. The data file contains definitions only. The
finite facts about the data are proved by \texttt{decide}, which the kernel checks; nothing is taken on trust from the
certificate \leanfile{tools/data/aperiodic_d6_dim25.json} or from the generator
\leanfile{tools/gen_aperiodic_example.py} (run it with \texttt{--check} to confirm that the data file is current).

\end{itemize}

\paragraph{Computations.} Every script that was used for this paper is in \leanfile{tools/}, and the saved codes are
in \leanfile{tools/data/}. The file \leanfile{tools/README.md} lists the scripts with their commands and the versions
used, and records the solver runs behind \cref{tab:small}, including those that timed out. Each script documents its
usage at its top.
\begin{itemize}[leftmargin=1.5em]\raggedright
\item \leanfile{tools/exhaustive_flag_checks.py} and \leanfile{tools/exhaustive_pigeonhole_checks.py} (Python with
NumPy): the exhaustive and randomized checks behind the numerical remarks of \cref{sec:order,sec:symmetry} and the
entries for $d\le4$ of \cref{tab:small}.
\item \leanfile{tools/three_term_example.py}: the examples of \cref{rem:three}.
\item \leanfile{tools/digit_and_composition_checks.py}: small instances of \cref{cor:digit} and of
\cref{rem:compose-TP}, both of which are proved in the text.
\item \leanfile{tools/entry_form_check.py}: small instances of \cref{prop:entry}, which is proved in Lean apart from
a count.
\item \leanfile{tools/two_class_family.py}: the rule of \cref{prop:two-class}, which is proved in the text, over five
small fields, with every rank stated in the proof; the saved codes over $\F_2$; and small instances of
\cref{cor:every-M} (\cref{rem:two-class}).
\item \leanfile{tools/class_uniform_sat.py} (needs the package \texttt{pycryptosat}): one run of the solver, for a
partition and a target dimension.
\item \leanfile{tools/small_sizes.py}: \cref{tab:small} with its bounds, recomputed, and with every saved code checked
again; the list of partitions up to affine maps; the search that raises the target dimension, which produced the saved
codes of the table; and controls that repeat the answers of the solver quoted in \cref{rem:small}.
\item \leanfile{tools/verify_class_uniform_code.py}: checks every saved code in \leanfile{tools/data/} against the
definition.
\item \leanfile{tools/analyse_code.py}: the intersections of the code of \cref{prop:aperiodic} with its translates,
quoted after that proposition.
\end{itemize}
Apart from the certificate of \cref{prop:aperiodic}, which Lean checks, these computations are not part of the formal
development, and no proof depends on them.


\bibliographystyle{plainurl}
\bibliography{references}

\end{document}
